\NormalFont\ShortTitle{Marcu}

{\MTB Vestea cea bună conform

\par }{\MT Marcu


\par }\ChapOne{1}{\PP \VerseOne{1}Începutul Veștii Bune a lui Iisus Hristos, Fiul lui Dumnezeu.


\par }{\PP \VS{2}După cum este scris în profeți,

\par }{\Q “Iată, trimit pe mesagerul Meu înaintea feței tale,

\par }{\QB care îți va pregăti calea înaintea ta:


\par }{\Q \VS{3}glasul cuiva care strigă în pustiu,

\par }{\QB 'Pregătiți calea Domnului!

\par }{\QB Fă-i cărările drepte!”


\par }{\PP \VS{4}Ioan a venit să boteze în pustie și să propovăduiască botezul pocăinței pentru iertarea păcatelor.

\VS{5}Tot ținutul Iudeii și toți cei din Ierusalim au ieșit la el. Au fost botezați de el în râul Iordan, mărturisindu-și păcatele.

\VS{6}Ioan era îmbrăcat cu păr de cămilă și avea o curea de piele la brâu. El mânca lăcuste și miere sălbatică.

\VS{7}El propovăduia, spunând: “După mine vine cel care este mai puternic decât mine, a cărui șnur al sandalelor nu sunt vrednic să mă aplec și să îl desfac.

\VS{8}Eu v-am botezat în apă, dar el vă va boteza în Duhul Sfânt.”


\par }{\PP \VS{9}În zilele acelea, Isus a venit din Nazaret din Galileea și a fost botezat de Ioan în Iordan.

\VS{10}Imediat ce a ieșit din apă, a văzut cerurile despărțindu-se și Duhul Sfânt coborând peste el ca un porumbel.

\VS{11}Un glas a venit din cer: “Tu ești Fiul Meu preaiubit, în care Îmi găsesc plăcerea.”


\par }{\PP \VS{12}Și îndată Duhul Sfânt l-a izgonit în pustiu.

\VS{13}Acolo, în pustiu, a stat patruzeci de zile, ispitit de Satana. Era cu animalele sălbatice; și îngerii îi slujeau.


\par }{\PP \VS{14}După ce Ioan a fost luat în temniță, Isus a venit în Galileea, propovăduind vestea cea bună a
{\WJ{Împărăției lui Dumnezeu,
}}

\VS{15}și zicând: “{\WJ{Timpul s-a împlinit și Împărăția lui Dumnezeu este aproape! Pocăiți-vă și credeți în Vestea cea Bună!”.
}}


\par }{\PP \VS{16}Trecând pe malul mării Galileii, a văzut pe Simon și pe Andrei, fratele lui Simon, care aruncau o plasă în mare, căci erau pescari.

\VS{17}Isus le-a zis:
{\WJ{“Veniți după Mine și vă voi face pescari de oameni”.
}}


\par }{\PP \VS{18}Și îndată și-au lăsat plasele și au mers după El.


\par }{\PP \VS{19}Și, mergând puțin mai departe de acolo, a văzut pe Iacov, fiul lui Zebedei, și pe fratele său Ioan, care erau și ei în corabie și reparau plasele.

\VS{20}Imediat i-a chemat, iar ei au lăsat pe tatăl lor, Zebedei, în barcă, cu slujitorii angajați, și au mers după el.


\par }{\PP \VS{21}S-au dus în Capernaum, și îndată, în ziua de Sabat, a intrat în sinagogă și învăța.

\VS{22}Ei erau uimiți de învățătura lui, pentru că îi învăța ca având autoritate, și nu ca pe cărturari.

\VS{23}Îndată s-a găsit în sinagoga lor un om cu un duh necurat și a strigat, zicând:

\VS{24}“Ha! Ce avem noi de-a face cu tine, Isus, Nazarineanule? Ai venit să ne distrugi? Eu știu cine ești: Sfântul lui Dumnezeu!”.


\par }{\PP \VS{25}Isus l-a mustrat și i-a zis: “{\WJ{Taci și ieși din el!”
}}


\par }{\PP \VS{26}Și duhul necurat a ieșit din el, care-l zguduia și striga cu glas tare.

\VS{27}Toți erau uimiți, așa că se întrebau între ei și ziceau: “Ce este aceasta? O învățătură nouă? Pentru că el poruncește cu autoritate chiar și duhurilor necurate, iar acestea îl ascultă!”

\VS{28}Imediat, vestea despre el s-a răspândit pretutindeni în toată regiunea Galileii și în împrejurimile ei.


\par }{\PP \VS{29}Și îndată ce au ieșit din sinagogă, au intrat în casa lui Simon și Andrei, cu Iacov și Ioan.

\VS{30}Mama soției lui Simon zăcea bolnavă de febră și îndată i-au povestit despre ea.

\VS{31}El a venit, a luat-o de mână și a ridicat-o. Febra a lăsat-o imediat și ea le-a slujit.


\par }{\PP \VS{32}Seara, când a apus soarele, au adus la El pe toți cei bolnavi și pe cei posedați de demoni.

\VS{33}Toată cetatea era adunată la ușă.

\VS{34}El a vindecat pe mulți bolnavi de diferite boli și a izgonit mulți demoni. Nu le permitea demonilor să vorbească, pentru că îl cunoșteau.


\par }{\PP \VS{35}Dis-de-dimineață, pe când era încă întuneric, s-a sculat și a ieșit afară, s-a dus într-un loc pustiu și s-a rugat acolo.

\VS{36}Simon și cei care erau cu el l-au căutat.

\VS{37}L-au găsit și i-au spus: “Toată lumea te caută.”


\par }{\PP \VS{38}Și le-a zis: “{\WJ{Să mergem în altă parte, în cetățile vecine, ca să propovăduiesc și acolo, căci pentru aceasta am ieșit.”
}}

\VS{39}Și a mers în sinagogile lor în toată Galileea, predicând și scoțând demoni.


\par }{\PP \VS{40}Un leprosar a venit la El și L-a rugat, îngenunchind în fața Lui și spunându-I: “Dacă vrei, poți să mă cureți.”


\par }{\PP \VS{41}Și, cuprins de milă, a întins mâna, l-a atins și i-a zis:
{\WJ{“Vreau să... Să fie curățat”.
}}

\VS{42}După ce a spus aceasta, îndată lepra s-a depărtat de la el și a fost curățit.

\VS{43}El l-a avertizat cu strictețe și l-a trimis imediat afară,

\VS{44}și i-a zis:
{\WJ{“Vezi să nu spui nimic nimănui, ci du-te să te prezinți preotului și oferă pentru curățirea ta ceea ce a poruncit Moise, ca mărturie pentru ei.”
}}


\par }{\PP \VS{45}Dar el a ieșit și a început să propovăduiască mult și să răspândească vestea, încât Isus nu mai putea intra în cetate, ci era afară, în locuri pustii. Oamenii veneau la el de pretutindeni.



\par }\Chap{2}{\PP \VerseOne{1}După câteva zile, când a intrat din nou în Capernaum, s-a auzit că este acasă.

\VS{2}Îndată s-au adunat mulți, încât nu mai era loc nici măcar în jurul ușii; și El le-a spus cuvântul.

\VS{3}Au venit patru oameni care îi duceau un paralitic.

\VS{4}Cum nu se puteau apropia de el din cauza mulțimii, au îndepărtat acoperișul în care se afla. După ce l-au spart, au dat drumul la covorul pe care era întins paraliticul.

\VS{5}Isus, văzând credința lor, i-a spus paraliticului:
{\WJ{“Fiule, păcatele tale îți sunt iertate.”
}}


\par }{\PP \VS{6}Dar unii din cărturari, care stăteau acolo, gândeau în inimile lor:

\VS{7}“Pentru ce vorbește omul acesta astfel de blasfemii? Cine poate ierta păcatele, dacă nu numai Dumnezeu?”


\par }{\PP \VS{8}Îndată Isus, văzând în duhul Său că ei gândeau astfel în ei înșiși, le-a zis: “{\WJ{Pentru ce gândiți astfel în inimile voastre?
}}

\VS{9}{\WJ{Ce este mai ușor, să-i spui
}}paraliticului:
{\WJ{“Păcatele tale sunt iertate”, sau să-i spui: “Scoală-te, ia-ți patul și umblă?”
}}

\VS{10}{\WJ{Dar ca să știți că Fiul Omului are putere pe pământ să ierte păcatele}}” — i-a spus paraliticului —

\VS{11}{\WJ{“Îți spun: “Scoală-te, ia-ți patul și du-te la casa ta”.
}}


\par }{\PP \VS{12}El s-a sculat, a luat îndată rogojina și a ieșit în fața tuturor, așa că toți au rămas uimiți și slăveau pe Dumnezeu, zicând: “Niciodată n-am văzut așa ceva.”


\par }{\PP \VS{13}Și a ieșit din nou la malul mării. Toată mulțimea a venit la el și el îi învăța.

\VS{14}Pe când trecea pe acolo, a văzut pe Levi, fiul lui Alfeu, care ședea la biroul fiscal. I-a spus:
{\WJ{“Vino după mine”. Și el s-a ridicat și l-a urmat.
}}


\par }{\PP \VS{15}El stătea la masă în casa lui, și mulți vameși și păcătoși s-au așezat la masă cu Isus și cu ucenicii Lui, căci erau mulți, și Îl urmau.

\VS{16}Cărturarii și fariseii, când au văzut că mănâncă cu păcătoșii și vameșii, au zis ucenicilor Lui: “De ce mănâncă și bea cu vameșii și cu păcătoșii?”


\par }{\PP \VS{17}Când a auzit Isus, le-a zis: “{\WJ{Cei sănătoși nu au nevoie de doctor, ci cei bolnavi. Eu nu am venit să chem pe cei drepți, ci pe păcătoși la pocăință.”
}}


\par }{\PP \VS{18}Ucenicii lui Ioan și fariseii posteau și au venit și L-au întrebat: “De ce postesc ucenicii lui Ioan și ucenicii fariseilor, iar ucenicii Tăi nu postesc?”


\par }{\PP \VS{19}Isus le-a zis: “{\WJ{Oare pot posti domnii de onoare, când mirele este cu ei? Atâta timp cât au mirele cu ei, nu pot posti.
}}

\VS{20}{\WJ{Dar vor veni zile când mirele va fi luat de la ei și atunci vor posti în ziua aceea.
}}

\VS{21}Nimeni nu coase
{\WJ{o bucată de pânză neîncălzită pe o haină veche, căci altfel peticul se micșorează și noul se desprinde de vechi și se face o gaură și mai mare.
}}

\VS{22}Nimeni nu pune vin nou în burdufuri vechi, căci
{\WJ{altfel vinul nou sparge burdufurile și vinul se varsă, iar burdufurile se distrug; dar ei pun vin nou în burdufuri proaspete.”
}}


\par }{\PP \VS{23}În ziua de Sabat, El mergea prin lanurile de grâu; și ucenicii Lui au început, pe când mergeau, să smulgă spice spicele.

\VS{24}Fariseii I-au zis: “Iată, de ce fac ei ceea ce nu este îngăduit în ziua Sabatului?”


\par }{\PP \VS{25}El le-a zis: “{\WJ{N-ați citit niciodată ce a făcut David și cei ce erau cu el, când a avut nevoie și a fost flămând?
}}

\VS{26}{\WJ{cum a intrat în casa lui Dumnezeu pe vremea marelui preot Abiatar și a mâncat pâinea de spectacol, pe care nu este îngăduit să o mănânce decât preoții, și a dat și celor care erau cu el?”
}}


\par }{\PP \VS{27}El le-a zis: “{\WJ{Sabatul a fost făcut pentru om, iar nu omul pentru Sabat.
}}

\VS{28}{\WJ{De aceea, Fiul Omului este stăpân chiar și pe Sabat.”
}}



\par }\Chap{3}{\PP \VerseOne{1}A intrat iarăși în sinagogă și era acolo un om căruia i se uscase mâna.

\VS{2}Îl pândeau, ca să vadă dacă îl va vindeca în ziua de Sabat, ca să-l acuze.

\VS{3}El i-a zis omului a cărui mână era uscată:
{\WJ{“Ridică-te!”
}}

\VS{4}El le-a zis: “{\WJ{Este oare îngăduit în ziua Sabatului să faci bine sau să faci rău? Să salvezi o viață sau să ucizi?” Dar ei au rămas tăcuți.
}}

\VS{5}După ce s-a uitat cu mânie la ei, mâhnit de împietrirea inimilor lor, i-a zis omului:
{\WJ{“Întinde-ți mâna.” El a întins-o și mâna i-a fost redată la fel de sănătoasă ca și cealaltă.
}}

\VS{6}Fariseii au ieșit afară și îndată au uneltit cu Irodienii împotriva lui, cum să-l nimicească.


\par }{\PP \VS{7}Isus S-a retras cu ucenicii Săi la mare; și o mare mulțime L-a urmat din Galileea, din Iudeea,

\VS{8}din Ierusalim, din Idumeea, de dincolo de Iordan, și din împrejurimile Tirului și Sidonului. O mare mulțime, auzind ce lucruri mari făcea, a venit la el.

\VS{9}El le-a spus discipolilor săi ca o barcă mică să rămână lângă el, din cauza mulțimii, ca să nu-l îmbrâncească.

\VS{10}Căci vindecase pe mulți, așa că toți cei care aveau boli se apăsau asupra lui ca să se atingă de el.

\VS{11}Spiritele necurate, de câte ori îl vedeau, cădeau înaintea lui și strigau: “Tu ești Fiul lui Dumnezeu!”

\VS{12}El i-a avertizat cu severitate să nu-l facă cunoscut.


\par }{\PP \VS{13}El s-a urcat pe munte și a chemat la Sine pe cei pe care-i voia și ei au venit la El.

\VS{14}A numit doisprezece, ca să fie cu el și să-i trimită să propovăduiască

\VS{15}și să aibă putere să vindece boli și să scoată demoni:

\VS{16}Simon (căruia i-a dat numele de Petru);

\VS{17}Iacov, fiul lui Zebedeu, și Ioan, fratele lui Iacov, (pe care i-a numit Boanerges, ceea ce înseamnă: “Fiii tunetului”);

\VS{18}Andrei, Filip, Bartolomeu, Matei, Toma, Iacov, fiul lui Alfeu, Tadeu, Simon Zelota,

\VS{19}și Iuda Iscarioteanul, care l-a și trădat.

\par }{\PP Apoi a intrat într-o casă.

\VS{20}Mulțimea s-a adunat din nou, încât nu mai putea nici măcar să mănânce pâine.

\VS{21}Când au auzit prietenii lui, au ieșit să-l prindă, căci ziceau: “Este nebun”.

\VS{22}Cărturarii care se coborau de la Ierusalim spuneau: “Are pe Beelzebul” și: “Prin prințul demonilor scoate afară demonii.”


\par }{\PP \VS{23}Apoi i-a chemat și le-a zis în pilde: “{\WJ{Cum poate Satana să scoată afară pe Satana?
}}

\VS{24}{\WJ{Dacă o împărăție este împărțită împotriva ei însăși, acea împărăție nu poate rămâne în picioare.
}}

\VS{25}{\WJ{Dacă o casă este împărțită împotriva ei însăși, casa aceea nu poate sta în picioare.
}}

\VS{26}{\WJ{Dacă Satana s-a ridicat împotriva lui însuși și este dezbinat, nu poate sta în picioare, ci are un sfârșit.
}}

\VS{27}{\WJ{Dar nimeni nu poate intra în casa celui puternic ca să jefuiască, dacă nu-l leagă mai întâi pe cel puternic; atunci îi va jefui casa.
}}


\par }{\PP \VS{28}“{\WJ{Adevărat vă spun că toate păcatele urmașilor omului vor fi iertate, inclusiv blasfemiile cu care vor huli;
}}

\VS{29}{\WJ{dar oricine va huli împotriva Duhului Sfânt nu va avea iertare, ci va fi supus osândei veșnice.”
}}

\VS{30}—pentru că au spus: “Are un duh necurat”.


\par }{\PP \VS{31}Mama și frații lui au venit și, stând afară, au trimis la el și l-au chemat.

\VS{32}O mulțime stătea în jurul lui și i-au spus: “Iată, mama ta, frații tăi și surorile tale sunt afară și te caută.”


\par }{\PP \VS{33}El le-a răspuns:
{\WJ{“Cine sunt mama Mea și frații Mei?”
}}

\VS{34}Și, uitându-se la cei care stăteau în jurul lui, a zis:
{\WJ{“Iată mama și frații mei!
}}

\VS{35}{\WJ{Căci oricine face voia lui Dumnezeu este fratele meu, sora mea și mama mea.”
}}



\par }\Chap{4}{\PP \VerseOne{1}Și a început din nou să învețe pe malul mării. S-a adunat la El o mare mulțime, așa că a intrat într-o corabie în mare și a șezut jos. Toată mulțimea era pe uscat, lângă mare.

\VS{2}El i-a învățat multe lucruri în parabole și le spunea în învățătura sa:

\VS{3}{\WJ{“Ascultați! Iată, agricultorul a ieșit să semene.
}}

\VS{4}În
{\WJ{timp ce semăna, o parte din sămânță a căzut pe drum, iar păsările au venit și au mâncat-o.
}}

\VS{5}{\WJ{Alte semințe au căzut pe un teren stâncos, unde era puțin pământ, și imediat au răsărit, pentru că nu aveau pământ adânc.
}}

\VS{6}{\WJ{Când a răsărit soarele, s-a pârjolit; și, pentru că nu avea rădăcină, s-a uscat.
}}

\VS{7}{\WJ{Alții au căzut între spini, iar spinii au crescut și l-au înăbușit, și nu a dat niciun rod.
}}

\VS{8}{\WJ{Altele au căzut în pământ bun și au dat rod, crescând și înmulțindu-se. Unii au produs de treizeci de ori, alții de șaizeci de ori, iar alții de o sută de ori mai mult.”
}}

\VS{9}El a spus:
{\WJ{“Cine are urechi de auzit, să audă.”
}}


\par }{\PP \VS{10}Când a rămas singur, cei ce erau în jurul lui cu cei doisprezece l-au întrebat despre pilde.

\VS{11}El le-a spus: “{\WJ{Vouă v-a fost dat misterul Împărăției lui Dumnezeu, dar pentru cei de afară, toate lucrurile se fac în parabole,
}}

\VS{12}pentru ca
{\WJ{“văzând să vadă și să nu priceapă, și auzind să audă și să nu înțeleagă, ca nu cumva să se întoarcă și să li se ierte păcatele”.”
}}


\par }{\PP \VS{13}El le-a zis: “{\WJ{Nu înțelegeți pilda aceasta? Cum veți înțelege toate parabolele?
}}

\VS{14}{\WJ{Fermierul seamănă cuvântul.
}}

\VS{15}Cei
{\WJ{de pe drum sunt cei în care este semănat cuvântul; și când au auzit, îndată vine Satana și ia cuvântul care a fost semănat în ei.
}}

\VS{16}{\WJ{Aceștia, la fel, sunt cei semănați pe locurile stâncoase, care, când au auzit cuvântul, îndată îl primesc cu bucurie.
}}

\VS{17}{\WJ{Ei nu au rădăcină în ei înșiși, ci sunt de scurtă durată. Când apare opresiunea sau persecuția din cauza cuvântului, imediat se poticnesc.
}}

\VS{18}{\WJ{Alții sunt cei care sunt semănați printre spini. Aceștia sunt cei care au ascultat cuvântul,
}}

\VS{19}dar
{\WJ{grijile acestui veac, înșelăciunea bogățiilor și poftele altor lucruri care intră în el înăbușă cuvântul și acesta devine neînsuflețit.
}}

\VS{20}{\WJ{Cei care au fost semănați pe pământ bun sunt cei care aud cuvântul, îl acceptă și fac rod, unii de treizeci de ori, alții de șaizeci de ori și alții de o sută de ori.”
}}


\par }{\PP \VS{21}El le-a zis: “{\WJ{Oare se aduce o lampă pentru a fi pusă sub coș sau sub pat? Nu este pusă pe un suport?
}}

\VS{22}{\WJ{Căci nimic nu este ascuns, decât ca să fie făcut cunoscut, și nimic nu a fost făcut secret, decât ca să iasă la lumină.
}}

\VS{23}{\WJ{Dacă are cineva urechi de auzit, să audă.”
}}


\par }{\PP \VS{24}Și le-a zis: “Luați aminte la
{\WJ{ce auziți. Cu orice măsură veți măsura, vi se va măsura și vouă; și vouă, celor ce ascultați, vi se va da mai mult.
}}

\VS{25}{\WJ{Căci oricui are, i se va da mai mult; iar celui care nu are, chiar și ceea ce are i se va lua.”
}}


\par }{\PP \VS{26}El a zis:
{\WJ{“Împărăția lui Dumnezeu este ca și cum un om ar arunca sămânță pe pământ,
}}

\VS{27}și ar
{\WJ{dormi și s-ar scula zi și noapte, iar sămânța ar răsări și ar crește, fără să știe cum.
}}

\VS{28}{\WJ{Pentru că pământul dă roade de la sine: mai întâi firul, apoi spicul, apoi bobul plin în spice.
}}

\VS{29}{\WJ{Dar când rodul este copt, îndată se pune secera, pentru că a venit secerișul.”
}}


\par }{\PP \VS{30}El a zis: “{\WJ{Cum vom asemăna noi Împărăția lui Dumnezeu? Sau cu ce parabolă o vom ilustra?
}}

\VS{31}{\WJ{Este ca un grăunte de muștar, care, deși este mai mic decât toate semințele care sunt pe pământ,
}}

\VS{32}{\WJ{totuși, după ce este semănat, crește și devine mai mare decât toate ierburile și scoate ramuri mari, încât păsările cerului pot să se adăpostească la umbra lui.”
}}


\par }{\PP \VS{33}Cu multe pilde ca acestea le spunea cuvântul, după cum puteau ei să-l audă.

\VS{34}Fără o pildă nu le vorbea; dar în particular, discipolilor Săi le explica totul.


\par }{\PP \VS{35}În ziua aceea, când s-a făcut seară, le-a zis:
{\WJ{“Să trecem dincolo.”
}}

\VS{36}Lăsând mulțimea, l-au luat cu ei, așa cum era, în barcă. Și alte bărci mici erau cu el.

\VS{37}S-a iscat o furtună mare de vânt și valurile băteau în barcă, atât de mult încât barca era deja plină.

\VS{38}El însuși era în pupa, dormind pe pernă; l-au trezit și l-au întrebat: “Învățătorule, nu-ți pasă că suntem pe moarte?”


\par }{\PP \VS{39}El s-a trezit, a mustrat vântul și a zis mării:
{\WJ{“Pace! Liniștește-te!” Vântul a încetat și a fost o mare acalmie.
}}

\VS{40}El le-a zis:
{\WJ{“De ce vă este atât de frică? Cum se face că nu aveți credință?”.
}}


\par }{\PP \VS{41}Ei s-au înspăimântat foarte tare și ziceau unul către altul: “Cine este Acesta, încât și vântul și marea Îl ascultă?”



\par }\Chap{5}{\PP \VerseOne{1}Și au ajuns de cealaltă parte a mării, în ținutul Gadarenilor.

\VS{2}După ce a ieșit din corabie, îndată i-a ieșit în întâmpinare un om cu un duh necurat din morminte.

\VS{3}El locuia în morminte. Nimeni nu-l mai putea lega, nici măcar cu lanțuri,

\VS{4}pentru că fusese de multe ori legat cu lanțuri și lanțuri, iar lanțurile fuseseră rupte de el și lanțurile rupte în bucăți. Nimeni nu mai avea puterea să-l îmblânzească.

\VS{5}Întotdeauna, zi și noapte, în morminte și în munți, striga și se tăia cu pietre.

\VS{6}Când l-a văzut de departe pe Isus, a alergat și s-a plecat în fața lui,

\VS{7}și, strigând cu glas tare, a zis: “Ce am eu de-a face cu tine, Isus, Fiul Dumnezeului Celui Preaînalt? Te conjur pe Dumnezeu, nu mă chinui!”.

\VS{8}Căci El i-a zis:
{\WJ{“Ieși din om, duh necurat!”
}}


\par }{\PP \VS{9}Și l-a întrebat:
{\WJ{“Cum te cheamă?”
}}

\par }{\PP El i-a spus: “Numele meu este Legiune, pentru că suntem mulți.”

\VS{10}Și l-a rugat mult ca să nu-i trimită afară din țară.

\VS{11}Or, pe coasta muntelui era o turmă mare de porci care se hrănea.

\VS{12}Toți demonii îl implorau, zicând: “Trimite-ne la porci, ca să intrăm în ei”.


\par }{\PP \VS{13}Atunci Isus le-a dat voie. Duhurile necurate au ieșit și au intrat în porci. Turma de aproximativ două mii de porci s-a repezit pe malul abrupt în mare și au fost înecați în mare.

\VS{14}Cei care hrăneau porcii au fugit și au povestit în oraș și în țară.

\par }{\PP Oamenii au venit să vadă ce s-a întâmplat.

\VS{15}Au venit la Isus și au văzut pe cel care fusese posedat de demoni stând jos, îmbrăcat și în deplinătatea facultăților mintale, pe cel care avea legionul; și s-au temut.

\VS{16}Cei care l-au văzut le-au povestit ce s-a întâmplat cu cel care era posedat de demoni și despre porci.

\VS{17}Ei au început să-l roage să plece din regiunea lor.


\par }{\PP \VS{18}Pe când se urca în corabie, cel ce fusese posedat de demoni l-a rugat să fie cu el.

\VS{19}El nu i-a îngăduit, ci i-a zis:
{\WJ{“Du-te la casa ta, la prietenii tăi, și spune-le ce lucruri mari a făcut Domnul pentru tine și cum a avut milă de tine.”
}}


\par }{\PP \VS{20}El a plecat și a început să vestească în Decapole cum a făcut Isus lucruri mari pentru el, și toată lumea se mira.


\par }{\PP \VS{21}După ce a trecut Isus cu corabia pe partea cealaltă, s-a adunat la El o mare mulțime de oameni, și El era pe malul mării.

\VS{22}Iată că a venit unul dintre conducătorii sinagogii, pe nume Iair; și, văzându-L, a căzut la picioarele Lui

\VS{23}și L-a rugat mult, zicând: “Fetița mea este în pragul morții. Te rog, vino și pune-ți mâinile peste ea, ca să se facă sănătoasă și să trăiască”.


\par }{\PP \VS{24}El a mers cu El și o mare mulțime de oameni Îl urmau și Îl strângeau din toate părțile.

\VS{25}O femeie care avea o scurgere de sânge de doisprezece ani,

\VS{26}și care suferise multe de la mulți doctori, cheltuise tot ce avea și nu se simțea mai bine, ci mai mult se înrăutățea,

\VS{27}și care, după ce a auzit cele despre Isus, s-a apropiat de El în mulțime și s-a atins de hainele Lui.

\VS{28}Căci ea zicea: “Dacă mă voi atinge de hainele Lui, mă voi face bine.”

\VS{29}Imediat, curgerea sângelui ei a secat și ea a simțit în trupul ei că este vindecată de suferința ei.


\par }{\PP \VS{30}Îndată Isus, văzând că puterea a ieșit din El, s-a întors în mulțime și a întrebat:
{\WJ{“Cine s-a atins de hainele Mele?”
}}


\par }{\PP \VS{31}Ucenicii Lui I-au zis: “Vezi mulțimea care Te strânge și zici: “Cine S-a atins de Mine?””


\par }{\PP \VS{32}S-a uitat în jur ca să o vadă pe cea care făcuse acest lucru.

\VS{33}Dar femeia, temându-se și tremurând, știind ce i se făcuse, a venit, a căzut înaintea lui și i-a spus tot adevărul.


\par }{\PP \VS{34}El i-a zis: “{\WJ{Fiică, credința ta te-a făcut sănătoasă. Du-te în pace și vindecă-te de boala ta.”
}}


\par }{\PP \VS{35}Pe când vorbea El încă, au venit niște oameni de la casa conducătorului sinagogii și au zis: “Fiica ta a murit. De ce să-l mai deranjezi pe Învățător?”


\par }{\PP \VS{36}Dar Isus, când a auzit mesajul rostit, a zis îndată conducătorului sinagogii:
{\WJ{“Nu te teme, ci crede!”
}}

\VS{37}Și nu a lăsat pe nimeni să-L urmeze, decât pe Petru, Iacov și Ioan, fratele lui Iacov.

\VS{38}A ajuns la casa conducătorului sinagogii și a văzut zarvă, plânsete și jale mare.

\VS{39}După ce a intrat înăuntru, le-a zis:
{\WJ{“De ce faceți zarvă și plângeți? Copilul nu este mort, ci a adormit”.
}}


\par }{\PP \VS{40}L-au batjocorit. Dar el, după ce i-a scos pe toți afară, a luat pe tatăl copilei, pe mama ei și pe cei care erau cu el și a intrat în locul unde zăcea copila.

\VS{41}Luând copila de mână, i-a zis:
{\WJ{“Talitha cumi!”,}} ceea ce înseamnă, interpretându-se: “{\WJ{Fetiță, îți spun, ridică-te!”.}}

\VS{42}Imediat fata s-a ridicat și a mers, căci avea doisprezece ani. Ei au rămas uimiți cu mare uimire.

\VS{43}El le-a poruncit cu strictețe ca nimeni să nu știe acest lucru și a poruncit să i se dea ceva de mâncare.



\par }\Chap{6}{\PP \VerseOne{1}Și a ieșit de acolo. A venit în țara lui și discipolii lui L-au urmat.

\VS{2}Când a venit Sabatul, a început să învețe în sinagogă și mulți care îl auzeau erau uimiți și ziceau: “De unde are omul acesta aceste lucruri?” și: “Ce înțelepciune i-a fost dată acestui om, de se fac prin mâinile lui asemenea lucruri mărețe?

\VS{3}Nu cumva acesta este tâmplarul, fiul Mariei și fratele lui Iacov, al lui Iose, al lui Iuda și al lui Simon? Nu cumva surorile lui sunt aici cu noi?” Așa că s-au supărat pe el.


\par }{\PP \VS{4}Isus le-a zis:
{\WJ{“Un prooroc nu este lipsit de cinste decât în țara lui, între rudele lui și în casa lui.”
}}

\VS{5}Acolo nu a putut face nicio lucrare măreață, decât că și-a pus mâinile peste câțiva bolnavi și i-a vindecat.

\VS{6}El se mira din cauza necredinței lor.

\par }{\PP A mers prin sate și a învățat.

\VS{7}A chemat la Sine pe cei doisprezece și a început să-i trimită doi câte doi; și le-a dat putere asupra duhurilor necurate.

\VS{8}Le-a poruncit să nu-și ia nimic pentru călătorie, în afară de un toiag, doar un toiag: nici pâine, nici portofel, nici bani în pungă,

\VS{9}ci să poarte sandale și să nu-și pună două tunici.

\VS{10}El le-a spus:
{\WJ{“Oriunde intrați într-o casă, rămâneți acolo până când veți pleca de acolo.
}}

\VS{11}{\WJ{Oricine nu vă va primi și nu vă va asculta, când plecați de acolo, scuturați praful de sub picioarele voastre ca mărturie împotriva lor. Cu siguranță, vă spun că, în ziua judecății, va fi mai ușor de suportat pentru Sodoma și Gomora decât pentru acest oraș!”
}}


\par }{\PP \VS{12}Ei au ieșit și propovăduiau oamenilor să se pocăiască.

\VS{13}Au scos mulți demoni, au uns cu untdelemn pe mulți bolnavi și i-au vindecat.

\VS{14}Regele Irod a auzit aceasta, căci numele lui devenise cunoscut, și a zis: “Ioan Botezătorul a înviat din morți și de aceea aceste puteri lucrează în el.”

\VS{15}Alții însă spuneau: “El este Ilie”. Alții au spus: “Este un profet sau ca unul dintre profeți”.

\VS{16}Dar Irod, când a auzit acestea, a zis: “Acesta este Ioan, pe care l-am decapitat. A înviat din morți”.

\VS{17}Căci însuși Irod trimisese să-l aresteze pe Ioan și-l legase în închisoare din cauza Irodiadei, soția fratelui său Filip, pentru că se căsătorise cu ea.

\VS{18}Căci Ioan îi spusese lui Irod: “Nu-ți este îngăduit să ai pe soția fratelui tău”.

\VS{19}Irodiada s-a pus împotriva lui și a vrut să-l ucidă, dar nu a putut,

\VS{20}pentru că Irod se temea de Ioan, știind că este un om drept și sfânt, și l-a păstrat în siguranță. Când îl auzea, făcea multe lucruri și îl asculta cu plăcere.


\par }{\PP \VS{21}Atunci a venit o zi potrivită, când Irod, de ziua lui de naștere, a dat o cină pentru nobilii săi, pentru arhierei și pentru căpeteniile Galileii.

\VS{22}Când însăși fiica Irodiadei a intrat și a dansat, a plăcut lui Irod și celor care stăteau cu el. Regele i-a spus tinerei: “Cere-mi tot ce vrei și îți voi da”.

\VS{23}El i-a jurat: “Orice îmi vei cere, îți voi da, până la jumătate din regatul meu.”


\par }{\PP \VS{24}Ea a ieșit și a zis mamei sale: “Ce să cer?”

\par }{\PP Ea a spus: “Capul lui Ioan Botezătorul”.


\par }{\PP \VS{25}Ea a venit îndată la rege și i-a zis: “Vreau să-mi dai acum capul lui Ioan Botezătorul pe un platou.”


\par }{\PP \VS{26}Împăratul a fost foarte mâhnit, dar, de dragul jurămintelor sale și al invitaților la masă, n-a vrut să o refuze.

\VS{27}Imediat, regele a trimis un soldat din garda sa și a poruncit să aducă capul lui Ioan; acesta s-a dus și l-a decapitat în închisoare,

\VS{28}i-a adus capul pe un platou și l-a dat tinerei domnișoare, iar tânăra domnișoară l-a dat mamei sale.


\par }{\PP \VS{29}Ucenicii Lui, auzind aceasta, au venit, au luat trupul Lui și l-au pus într-un mormânt.


\par }{\PP \VS{30}Apostolii s-au adunat la Isus și I-au povestit tot ce făcuseră și tot ce învățaseră.

\VS{31}El le-a zis:
{\WJ{“Veniți într-un loc pustiu și odihniți-vă puțin.” Căci erau mulți care veneau și plecau, și nu aveau timp liber nici măcar să mănânce.
}}

\VS{32}Ei s-au dus cu barca într-un loc pustiu, singuri.

\VS{33}Văzându-i mergând, mulți l-au recunoscut și au alergat acolo pe jos din toate cetățile. Au ajuns înaintea lor și au venit împreună la el.

\VS{34}Isus a ieșit, a văzut o mulțime mare și i s-a făcut milă de ei, pentru că erau ca niște oi fără păstor; și a început să-i învețe multe lucruri.

\VS{35}Când s-a făcut târziu, discipolii Lui au venit la El și i-au spus: “Locul acesta este pustiu și este târziu.

\VS{36}Trimiteți-i afară, ca să meargă în ținutul și în satele din jur și să-și cumpere pâine, căci nu au ce mânca.”


\par }{\PP \VS{37}Dar El le-a răspuns: “{\WJ{Dați-le voi ceva de mâncare.”
}}

\par }{\PP L-au întrebat: “Să mergem să cumpărăm pâine în valoare de două sute de denari și să le dăm ceva de mâncare?”.


\par }{\PP \VS{38}El le-a zis: “{\WJ{Câte pâini aveți? Mergeți să vedeți.”
}}

\par }{\PP Când au aflat, au spus: “Cinci și doi pești”.


\par }{\PP \VS{39}Și le-a poruncit să se așeze toți în grupuri pe iarba verde.

\VS{40}Și s-au așezat în rânduri, cu sutele și cu cincizeci.

\VS{41}A luat cele cinci pâini și cei doi pești; și, privind spre cer, a binecuvântat și a frânt pâinile, pe care le-a dat ucenicilor Săi ca să le pună înaintea lor, iar cei doi pești i-a împărțit la toți.

\VS{42}Toți au mâncat și s-au săturat.

\VS{43}Au luat douăsprezece coșuri pline cu bucățile frânte și cu peștii.

\VS{44}Cei care au mâncat pâinile au fost cinci mii de oameni.


\par }{\PP \VS{45}Și îndată a pus pe ucenicii Săi să se urce în corabie și să meargă mai departe, în cealaltă parte, la Betsaida, în timp ce El însuși a trimis mulțimea.

\VS{46}După ce și-a luat rămas bun de la ei, s-a urcat pe munte ca să se roage.


\par }{\PP \VS{47}Când s-a făcut seară, corabia era în mijlocul mării, iar el era singur pe uscat.

\VS{48}Văzându-i chinuiți la vâslit, pentru că vântul le era potrivnic, pe la a patra strajă a nopții, a venit la ei, mergând pe mare; și ar fi vrut să treacă pe lângă ei,

\VS{49}dar ei, când l-au văzut mergând pe mare, au crezut că este o fantomă și au strigat;

\VS{50}pentru că toți l-au văzut și s-au tulburat. Dar el a vorbit imediat cu ei și le-a spus:
{\WJ{“Înveseliți-vă! Eu sunt! Nu vă temeți!”.
}}

\VS{51}S-a urcat cu ei în barcă, iar vântul a încetat, iar ei erau foarte mirați între ei și se minunau;

\VS{52}pentru că nu înțeleseseră ce se întâmplase cu pâinile, dar aveau inima împietrită.


\par }{\PP \VS{53}După ce au trecut dincolo, au ajuns la țărm, la Genesaret, și au acostat la țărm.

\VS{54}După ce au ieșit din barcă, oamenii l-au recunoscut imediat,

\VS{55}și au alergat în toată regiunea aceea și au început să-i aducă pe cei bolnavi pe saltelele lor acolo unde auziseră că se află.

\VS{56}Oriunde intra — în sate, în orașe sau la țară —, ei îi așezau pe bolnavi în piețe și îl rugau să atingă doar franjurii hainei lui; și toți cei care se atingeau de el se făceau bine.



\par }\Chap{7}{\PP \VerseOne{1}Atunci s-au adunat la El fariseii și unii din cărturari, veniți de la Ierusalim.

\VS{2}Când au văzut pe unii dintre ucenicii Săi mâncând pâine cu mâinile spurcate, adică nespălate, au găsit o vină.

\VS{3}(Căci fariseii și toți iudeii nu mănâncă dacă nu-și spală mâinile și antebrațele, ținându-se de tradiția bătrânilor.

\VS{4}Ei nu mănâncă atunci când vin de la piață dacă nu se spală, și mai sunt multe alte lucruri pe care au primit să le țină: spălările paharelor, ale ulcioarelor, ale vaselor de bronz și ale canapelelor).

\VS{5}Fariseii și cărturarii l-au întrebat: “De ce nu umblă discipolii tăi după tradiția bătrânilor, ci mănâncă pâinea lor cu mâinile nespălate?”


\par }{\PP \VS{6}El le-a răspuns: “{\WJ{Bine a proorocit Isaia despre voi, fățarnicilor, cum este scris,
}}

\par }{\Q {\WJ{'Acest popor mă cinstește cu buzele lui,
}}

\par }{\QB {\WJ{dar inima lor este departe de mine.
}}


\par }{\Q \VS{7}{\WJ{În zadar se închină ei la Mine,
}}

\par }{\QB {\WJ{învățând ca doctrine poruncile oamenilor.
}}


\par }{\PP \VS{8}{\WJ{“Căci voi lăsați deoparte porunca lui Dumnezeu și vă țineți de tradiția oamenilor, spălând ulcioarele și paharele, și faceți multe alte lucruri de acest fel.”
}}

\VS{9}El le-a zis: “{\WJ{Foarte bine că respingeți porunca lui Dumnezeu, ca să vă păstrați tradiția voastră.
}}

\VS{10}{\WJ{Căci Moise a zis: “Cinstește pe tatăl tău și pe mama ta”; și: “Cine vorbește de rău pe tatăl sau pe mama, să fie omorât”.
}}

\VS{11}{\WJ{Dar voi ziceți: “Dacă un om spune tatălui său sau mamei sale: “Orice folos pe care l-ai fi putut primi de la mine este Corban””,
}}adică dat lui Dumnezeu,

\VS{12}{\WJ{“atunci nu-i mai permiteți să facă nimic pentru tatăl său sau pentru mama sa,
}}

\VS{13}{\WJ{anulând cuvântul lui Dumnezeu prin tradiția voastră pe care ați transmis-o”. Voi faceți multe lucruri de acest fel”.
}}


\par }{\PP \VS{14}A chemat toată mulțimea la Sine și le-a zis: “{\WJ{Ascultați-mă cu toții și înțelegeți.
}}

\VS{15}{\WJ{Nu este nimic din afara omului care, intrând în el, să îl spurce; dar lucrurile care ies din om sunt cele care îl spurcă pe om.
}}

\VS{16}{\WJ{Dacă are cineva urechi de auzit, să audă!”
}}


\par }{\PP \VS{17}După ce a intrat într-o casă, departe de mulțime, ucenicii Lui L-au întrebat despre pilda aceea.

\VS{18}El le-a zis: “Sunteți
{\WJ{și voi fără pricepere? Nu vă dați seama că tot ceea ce intră în om din afară nu-l poate spurca,
}}

\VS{19}{\WJ{pentru că nu intră în inima lui, ci în stomac, apoi în latrină, ceea ce face ca toate alimentele să fie curate?”
}}

\VS{20}El a răspuns: “Ceea ce
{\WJ{iese din om, aceea spurcă omul.
}}

\VS{21}{\WJ{Căci dinăuntru, din inima omului, ies gândurile rele, adulterele, păcatele sexuale, crimele, furturile,
}}

\VS{22}{\WJ{poftele, răutatea, înșelăciunea, poftele, ochiul rău, blasfemia, mândria și nebunia.
}}

\VS{23}{\WJ{Toate aceste lucruri rele vin dinăuntru și îl spurcă pe om.”
}}


\par }{\PP \VS{24}De acolo s-a sculat și a plecat în ținutul Tirului și al Sidonului. A intrat într-o casă și n-a vrut să știe nimeni, dar n-a putut scăpa de observație.

\VS{25}Căci o femeie a cărei fetiță avea un duh necurat, auzind de el, a venit și s-a aruncat la picioarele lui.

\VS{26}Femeia era o greacă, siro-feniciană de rasă. Ea l-a implorat să alunge demonul din fiica ei.

\VS{27}Dar Isus i-a zis:
{\WJ{“Lasă să se sature mai întâi copiii, căci nu se cuvine să iei pâinea copiilor și să o arunci la câini.”
}}


\par }{\PP \VS{28}Dar ea i-a răspuns: “Da, Doamne. Totuși, chiar și câinii de sub masă mănâncă firimiturile copiilor.”


\par }{\PP \VS{29}El i-a zis: “{\WJ{Pentru această vorbă, du-te. Demonul a ieșit din fiica ta.”
}}


\par }{\PP \VS{30}Ea s-a dus la casa ei și a găsit copilul culcat pe pat, iar demonul ieșise.


\par }{\PP \VS{31}A plecat din nou de la hotarele Tirului și ale Sidonului și a ajuns la marea Galileii, trecând prin mijlocul ținutului Decapolei.

\VS{32}I-au adus pe unul care era surd și avea o împiedicare în vorbire. L-au rugat să pună mâna pe el.

\VS{33}L-a luat la o parte de la mulțime, în particular, și i-a băgat degetele în urechi; apoi a scuipat și i-a atins limba.

\VS{34}Privind spre cer, a suspinat și i-a zis:
{\WJ{“Ephatha!”,}} adică:
{\WJ{“Deschide-te!”.}}

\VS{35}Imediat i s-au deschis urechile, i s-a eliberat impedimentul limbii și a vorbit limpede.

\VS{36}El le-a poruncit să nu spună nimănui, dar cu cât le poruncea mai mult, cu atât mai mult o vesteau.

\VS{37}Ei au fost uimiți peste măsură, spunând: “A făcut toate lucrurile bine. El îi face chiar și pe surzi să audă și pe cei muți să vorbească!”



\par }\Chap{8}{\PP \VerseOne{1}În zilele acelea, pe când era o mulțime foarte mare și nu aveau ce mânca, Isus a chemat la Sine pe ucenicii Săi și le-a zis:

\VS{2}“Mi-e
{\WJ{milă de mulțime, pentru că de trei zile stau la Mine și nu au ce mânca.
}}

\VS{3}{\WJ{Dacă îi voi trimite cu post la casele lor, vor leșina pe drum, pentru că unii dintre ei au venit de departe.”
}}


\par }{\PP \VS{4}Ucenicii Lui I-au răspuns: “De unde ar putea cineva să sature pe acești oameni cu pâine aici, într-un loc pustiu?”


\par }{\PP \VS{5}Și i-a întrebat: “{\WJ{Câte pâini aveți?”
}}

\par }{\PP Ei au spus: “Șapte”.


\par }{\PP \VS{6}A poruncit mulțimii să șadă pe pământ și a luat cele șapte pâini. După ce a mulțumit, le-a frânt și le-a dat ucenicilor Săi să le servească, iar ei au servit mulțimii.

\VS{7}Au luat și câțiva peștișori. După ce i-a binecuvântat, a spus să servească și pe aceștia.

\VS{8}Ei au mâncat și s-au săturat. Au luat șapte coșuri cu bucățile sparte care rămăseseră.

\VS{9}Cei care mâncaseră erau în jur de patru mii. Apoi i-a trimis departe.


\par }{\PP \VS{10}Și îndată a intrat în corabie cu ucenicii Săi și a ajuns în ținutul Dalmanutha.

\VS{11}Fariseii au ieșit și au început să-L întrebe, căutând de la El un semn din cer și punându-L la încercare.

\VS{12}El a suspinat adânc în duhul său și a zis:
{\WJ{“De ce caută neamul acesta un semn? Cu siguranță vă spun că nu se va da niciun semn acestei generații”.
}}


\par }{\PP \VS{13}I-a lăsat și, intrând din nou în corabie, a plecat în cealaltă parte.

\VS{14}Au uitat să ia pâine; și nu aveau cu ei în barcă mai mult de o pâine.

\VS{15}El i-a avertizat, zicând:
{\WJ{“Luați seama: păziți-vă de drojdia fariseilor și de drojdia lui Irod.”
}}


\par }{\PP \VS{16}Ei se certau între ei, zicând: “Pentru că nu avem pâine.”


\par }{\PP \VS{17}Isus, văzând aceasta, le-a zis: “{\WJ{De ce credeți că este din pricină că nu aveți pâine? Nu vă dați seama încă și nu înțelegeți? Inima voastră este încă împietrită?
}}

\VS{18}{\WJ{Având ochi, nu vedeți? Având urechi, nu auziți? Nu vă amintiți?
}}

\VS{19}{\WJ{Când am împărțit cele cinci pâini la cele cinci mii de oameni, câte coșuri pline de bucăți ați luat?”
}}

\par }{\PP I-au spus: “Douăsprezece”.


\par }{\PP \VS{20}{\WJ{“Când cele șapte pâini au hrănit cele patru mii de oameni, câte coșuri pline cu frânturi ați luat?”
}}

\par }{\PP I-au spus: “Șapte”.


\par }{\PP \VS{21}Și i-a întrebat: “{\WJ{Nu înțelegeți încă?”
}}


\par }{\PP \VS{22}Și a venit la Betsaida. I-au adus un orb și l-au rugat să se atingă de el.

\VS{23}El a apucat pe orb de mână și l-a scos din sat. După ce l-a scuipat pe ochi și și-a pus mâinile pe el, l-a întrebat dacă a văzut ceva.


\par }{\PP \VS{24}Și s-a uitat în sus și a zis: “Văd oameni, dar îi văd ca pe niște copaci care umblă.”


\par }{\PP \VS{25}Apoi și-a pus iarăși mâinile pe ochi. El s-a uitat cu atenție și, după ce s-a refăcut, a văzut clar pe toată lumea.

\VS{26}L-a trimis la casa lui, spunându-i:
{\WJ{“Să nu intri în sat și să nu spui nimănui din sat.”
}}


\par }{\PP \VS{27}Isus a ieșit cu ucenicii Săi în satele din Cezareea lui Filipi. Pe drum i-a întrebat pe discipolii Săi:
{\WJ{“Cine spun oamenii că sunt Eu?”.
}}


\par }{\PP \VS{28}Ei i-au spus: “Ioan Botezătorul; alții zic Ilie, iar alții zic Ilie, iar alții, unul din profeți.”


\par }{\PP \VS{29}El le-a zis: “{\WJ{Dar voi cine ziceți că sunt Eu?”
}}

\par }{\PP Petru i-a răspuns: “Tu ești Hristosul.”


\par }{\PP \VS{30}Și le-a poruncit să nu spună nimănui despre El.

\VS{31}A început să-i învețe că Fiul Omului trebuie să sufere multe, să fie lepădat de bătrâni, de preoții cei mai de seamă și de cărturari, să fie omorât și să învie după trei zile.

\VS{32}El le vorbea deschis. Petru l-a luat și a început să îl dojenească.

\VS{33}Dar el, întorcându-se și văzându-i pe ucenicii săi, l-a mustrat pe Petru și i-a zis: “{\WJ{Treci în spatele meu, Satana! Căci tu nu te gândești la lucrurile lui Dumnezeu, ci la lucrurile oamenilor”.
}}


\par }{\PP \VS{34}A chemat mulțimea la Sine împreună cu ucenicii Săi și le-a zis: “{\WJ{Cine vrea să vină după Mine, să se lepede de sine, să-și ia crucea și să-Mi urmeze Mie.
}}

\VS{35}{\WJ{Căci oricine vrea să-și salveze viața o va pierde, iar oricine își va pierde viața pentru Mine și pentru Buna Vestire o va salva.
}}

\VS{36}{\WJ{Căci ce-i folosește unui om să câștige toată lumea și să-și piardă viața?
}}

\VS{37}{\WJ{Căci ce va da omul în schimbul vieții sale?
}}

\VS{38}{\WJ{Căci oricine se va rușina de mine și de cuvintele mele în acest neam adulterin și păcătos, și Fiul Omului se va rușina de el când va veni în slava Tatălui său cu sfinții îngeri.”
}}



\par }\Chap{9}{\PP \VerseOne{1}El le-a zis: “{\WJ{Adevărat vă spun că sunt unii care stau aici și care nu vor gusta moartea până nu vor vedea venind cu putere Împărăția lui Dumnezeu.”
}}


\par }{\PP \VS{2}După șase zile, Isus a luat cu El pe Petru, Iacov și Ioan și i-a dus singuri pe un munte înalt, în particular, și S-a schimbat în fața lor.

\VS{3}Hainele Lui au devenit strălucitoare, extrem de albe, ca zăpada, așa cum nici un spălător de pe pământ nu le poate albi.

\VS{4}Le-au apărut Ilie și Moise, care stăteau de vorbă cu Isus.


\par }{\PP \VS{5}Petru a răspuns lui Isus: “Rabi, este bine să fim aici. Să facem trei corturi: unul pentru tine, unul pentru Moise și unul pentru Ilie.”

\VS{6}Căci nu știa ce să spună, pentru că erau foarte speriați.


\par }{\PP \VS{7}Un nor a venit și i-a acoperit cu umbră, iar din nor a ieșit un glas: “Acesta este Fiul Meu preaiubit. Ascultați-l.”


\par }{\PP \VS{8}Și, uitându-se deodată în jur, n-au mai văzut pe nimeni cu ei, decât numai pe Isus.


\par }{\PP \VS{9}Pe când se coborau de pe munte, le-a poruncit să nu spună nimănui cele ce văzuseră, până ce Fiul Omului nu va învia din morți.

\VS{10}Ei au păstrat pentru ei înșiși acest cuvânt, întrebându-se ce înseamnă “învierea din morți”.


\par }{\PP \VS{11}Și L-au întrebat: “Pentru ce zic cărturarii că Ilie trebuie să vină mai întâi?”


\par }{\PP \VS{12}{\WJ{El
}}le-a zis: “{\WJ{Ilie vine mai întâi și restaurează toate lucrurile. Cum este scris despre Fiul Omului că trebuie să sufere multe și să fie disprețuit?
}}

\VS{13}{\WJ{Dar Eu vă spun că Ilie a venit și i-au făcut și lui tot ce au vrut, așa cum este scris despre El.”
}}


\par }{\PP \VS{14}Când a venit la ucenici, a văzut o mulțime mare în jurul lor și cărturari care îi întrebau.

\VS{15}Îndată toată mulțimea, văzându-l, a fost foarte uimită și, alergând la el, l-a salutat.

\VS{16}El i-a întrebat pe cărturari:
{\WJ{“Ce-i întrebați?”
}}


\par }{\PP \VS{17}Unul din mulțime a răspuns: “Învățătorule, ți-am adus pe fiul meu, care are un duh mut,

\VS{18}și oriunde îl apucă, îl aruncă la pământ; face spume la gură, scrâșnește din dinți și se înțeapă. I-am rugat pe discipolii tăi să-l scoată afară, dar n-au fost în stare.”


\par }{\PP \VS{19}El i-a răspuns: “{\WJ{Neam necredincios, până când voi fi cu voi? Până când voi mai suporta cu voi? Aduceți-l la Mine!”
}}


\par }{\PP \VS{20}L-au adus la el și, când l-a văzut, îndată l-a cuprins duhul și a căzut la pământ, zvârcolindu-se și făcând spume la gură.


\par }{\PP \VS{21}Și a întrebat pe tatăl său: “De
{\WJ{când i se întâmplă asta?”
}}

\par }{\PP El a spus: “Din copilărie.

\VS{22}De multe ori l-a aruncat atât în foc, cât și în apă, pentru a-l distruge. Dar dacă poți face ceva, ai milă de noi și ajută-ne.”


\par }{\PP \VS{23}Isus i-a zis: “{\WJ{Dacă poți crede, toate sunt cu putință celui ce crede.”
}}


\par }{\PP \VS{24}Și îndată tatăl copilului a strigat cu lacrimi: “Cred. Ajută-mi necredința!”


\par }{\PP \VS{25}Isus, văzând că o mulțime alerga laolaltă, a mustrat pe duhul necurat și i-a zis: “{\WJ{Duh mut și surd, îți poruncesc să ieși din el și să nu mai intri niciodată în el!”
}}


\par }{\PP \VS{26}Și după ce a strigat și a fost cuprins de mari convulsii, a ieșit din el. Băiatul s-a făcut ca un mort, atât de mult încât cei mai mulți dintre ei au spus: “A murit”.

\VS{27}Dar Isus l-a luat de mână și l-a înviat; și s-a ridicat.


\par }{\PP \VS{28}După ce a intrat în casă, ucenicii Lui L-au întrebat în particular: “De ce n-am putut să-l scoatem afară?”


\par }{\PP \VS{29}Și le-a zis: “Nu se
{\WJ{poate ieși așa ceva decât prin rugăciune și prin post.”
}}


\par }{\PP \VS{30}Au plecat de acolo și au trecut prin Galileea. El nu voia să știe nimeni,

\VS{31}pentru că îi învăța pe ucenicii Săi și le spunea: “{\WJ{Fiul Omului va fi dat în mâinile oamenilor și ei Îl vor ucide; iar după ce va fi ucis, a treia zi va învia.”
}}


\par }{\PP \VS{32}Dar ei nu înțelegeau cuvântul acesta și se temeau să-L întrebe.


\par }{\PP \VS{33}Când a ajuns în Capernaum, i-a întrebat pe cei din casă: “{\WJ{Ce vă certați între voi pe drum?”
}}


\par }{\PP \VS{34}Dar ei au tăcut, pentru că pe drum se certau între ei cine este cel mai mare.


\par }{\PP \VS{35}A șezut jos, a chemat pe cei doisprezece și le-a zis: “{\WJ{Dacă cineva vrea să fie cel dintâi, să fie cel din urmă dintre toți și slujitorul tuturor.”
}}

\VS{36}A luat un copilaș și l-a așezat în mijlocul lor. Luându-l în brațe, le-a spus:

\VS{37}{\WJ{“Oricine primește un astfel de copilaș în numele meu, pe mine mă primește; și oricine mă primește pe mine, nu pe mine mă primește, ci pe cel care m-a trimis pe mine.”
}}


\par }{\PP \VS{38}Ioan I-a zis: “Învățătorule, am văzut pe cineva care nu ne urmează, care scoate demoni în Numele Tău, și i-am interzis, pentru că nu ne urmează.”


\par }{\PP \VS{39}Dar Isus a zis: “{\WJ{Nu-i interziceți, căci nu este nimeni care să facă o lucrare mare în Numele Meu și să poată repede să vorbească de rău despre Mine.
}}

\VS{40}{\WJ{Căci oricine nu este împotriva noastră este de partea noastră.
}}

\VS{41}{\WJ{Căci oricine vă va da să beți un pahar cu apă în numele meu, pentru că sunteți ai lui Hristos, cu siguranță vă spun că nu-și va pierde în niciun fel răsplata.
}}


\par }{\PP \VS{42}{\WJ{Oricine va face să se poticnească pe unul din acești mici care cred în Mine, mai bine ar fi să fie aruncat în mare, cu o piatră de moară atârnată de gât.
}}

\VS{43}{\WJ{Dacă mâna ta te face să te poticnești, taie-ți-o! Este mai bine pentru tine să intri în viață mutilat, decât să ai cele două mâini pentru a merge în Gheena, în focul nestins,
}}

\VS{44}{\WJ{“unde viermele lor nu moare și focul nu se stinge”.
}}

\VS{45}{\WJ{Dacă piciorul tău te face să te poticnești, taie-l! Este mai bine pentru tine să intri în viață șchiop, decât ca cele două picioare să fie aruncate în Gheena, în focul care nu se va stinge niciodată —
}}

\VS{46}{\WJ{'unde viermele lor nu moare și focul nu se stinge'.
}}

\VS{47}{\WJ{Dacă ochiul tău te face să te poticnești, aruncă-l afară. Este mai bine pentru tine să intri în Împărăția lui Dumnezeu cu un singur ochi, decât să ai doi ochi pentru a fi aruncat în Gheena de foc,
}}

\VS{48}{\WJ{'unde viermele lor nu moare și focul nu se stinge'.
}}

\VS{49}{\WJ{Căci toată lumea va fi sărată cu foc și orice jertfă va fi condimentată cu sare.
}}

\VS{50}{\WJ{Sarea este bună, dar, dacă sarea și-a pierdut savoarea, cu ce o vei asezona? Aveți sare în voi înșivă și fiți în pace unii cu alții.”
}}



\par }\Chap{10}{\PP \VerseOne{1}S-a ridicat de acolo și a ajuns în ținutul Iudeii și dincolo de Iordan. Mulțimile s-au adunat din nou la El. Așa cum făcea de obicei, îi învăța din nou.


\par }{\PP \VS{2}Fariseii, venind la El, L-au pus la încercare și L-au întrebat: “Este îngăduit ca un bărbat să divorțeze de nevasta sa?”


\par }{\PP \VS{3}El a răspuns:
{\WJ{“Ce ți-a poruncit Moise?”
}}


\par }{\PP \VS{4}Ei au zis: “Moise a îngăduit să se scrie un certificat de divorț și să se divorțeze de ea.”


\par }{\PP \VS{5}Dar Isus le-a zis: “{\WJ{Pentru împietrirea inimii voastre, v-a scris porunca aceasta.
}}

\VS{6}{\WJ{Dar, de la începutul creației, Dumnezeu i-a făcut bărbat și femeie.
}}

\VS{7}De
{\WJ{aceea, omul va lăsa pe tatăl său și pe mama sa și se va uni cu soția sa,
}}

\VS{8}{\WJ{și cei doi se vor face un singur trup, astfel încât nu mai sunt doi, ci un singur trup.
}}

\VS{9}{\WJ{Așadar, ceea ce Dumnezeu a unit, nimeni să nu despartă.”
}}


\par }{\PP \VS{10}În casă, ucenicii Lui L-au întrebat din nou despre același lucru.

\VS{11}El le-a spus: “{\WJ{Oricine divorțează de soția sa și se căsătorește cu alta comite adulter împotriva ei.
}}

\VS{12}{\WJ{Dacă o femeie însăși divorțează de soțul ei și se căsătorește cu altul, comite adulter.”
}}


\par }{\PP \VS{13}Îi aduceau prunci, ca să se atingă de ei; dar ucenicii mustrau pe cei ce-i aduceau.

\VS{14}Dar Isus, văzând aceasta, s-a indignat și le-a zis: “{\WJ{Lăsați copilașii să vină la Mine! Nu le interziceți, căci Împărăția lui Dumnezeu aparține unora ca aceștia.
}}

\VS{15}{\WJ{Adevărat vă spun că oricine nu va primi Împărăția lui Dumnezeu ca un copilaș, nu va intra nicidecum în ea.”
}}

\VS{16}El i-a luat în brațe și i-a binecuvântat, punându-și mâinile peste ei.


\par }{\PP \VS{17}Pe când ieșea El pe cale, a alergat unul la El, a îngenuncheat înaintea Lui și L-a întrebat: “Învățătorule, ce să fac ca să moștenesc viața veșnică?”


\par }{\PP \VS{18}Isus i-a zis:
{\WJ{“De ce Mă numești bun? Nimeni nu este bun, în afară de unul singur: Dumnezeu.
}}

\VS{19}{\WJ{Tu cunoști poruncile: Să nu ucizi, să nu comiți adulter, să nu furi, să nu depui mărturie mincinoasă, să nu înșeli, să cinstești pe tatăl tău și pe mama ta.”
}}


\par }{\PP \VS{20}El I-a zis: “Învățătorule, toate acestea le știu din tinerețe.”


\par }{\PP \VS{21}Isus, uitându-se la el, l-a iubit și i-a zis: “{\WJ{Un singur lucru îți lipsește. Du-te, vinde tot ce ai și dă săracilor și vei avea o comoară în ceruri; și vino, urmează-Mi, luându-ți crucea.”
}}


\par }{\PP \VS{22}Dar, când a auzit acest cuvânt, fața lui a căzut și a plecat mâhnit, pentru că era un om care avea multe averi.


\par }{\PP \VS{23}Isus s-a uitat în jur și a zis ucenicilor Săi: “{\WJ{Cât de greu intră în Împărăția lui Dumnezeu cei ce au bogății!”
}}


\par }{\PP \VS{24}Ucenicii au rămas uimiți de cuvintele Lui. Dar Isus le-a răspuns din nou: “{\WJ{Copii, cât de greu intră în Împărăția lui Dumnezeu cei ce se încred în bogății!
}}

\VS{25}{\WJ{Este mai ușor pentru o cămilă să treacă prin urechea unui ac decât pentru un om bogat să intre în Împărăția lui Dumnezeu.”
}}


\par }{\PP \VS{26}Ei au fost foarte mirați și I-au zis: “Cine poate fi mântuit?”


\par }{\PP \VS{27}Isus, uitându-se la ei, a zis: “La
{\WJ{oameni este cu neputință, dar nu și la Dumnezeu, căci la Dumnezeu toate lucrurile sunt cu putință.”
}}


\par }{\PP \VS{28}Și Petru a început să-i spună: “Iată, noi am lăsat totul și am venit după Tine.”


\par }{\PP \VS{29}Isus a zis: “{\WJ{Adevărat vă spun că nu este nimeni care să fi lăsat casă, frați, surori, tată, mamă, nevastă, copii sau pământ, pentru Mine și pentru Buna Vestire,
}}

\VS{30}{\WJ{ci va primi de o sută de ori mai mult acum, în veacul acesta: case, frați, surori, mame, copii, pământ, cu prigoane, iar în veacul viitor, viața veșnică.
}}

\VS{31}{\WJ{Dar mulți dintre cei dintâi vor fi cei din urmă, iar cei din urmă cei dintâi.”
}}


\par }{\PP \VS{32}Ei erau pe drum, urcându-se la Ierusalim; și Isus mergea înaintea lor, și ei erau uimiți, iar cei ce-i urmau se temeau. El i-a luat din nou pe cei doisprezece și a început să le spună lucrurile care aveau să i se întâmple.

\VS{33}{\WJ{“Iată, ne suim la Ierusalim. Fiul Omului va fi dat pe mâna preoților de seamă și a cărturarilor. Aceștia îl vor condamna la moarte și îl vor da pe mâna neamurilor.
}}

\VS{34}Îl
{\WJ{vor batjocori, îl vor scuipa, îl vor biciui și îl vor ucide. A treia zi va învia.”
}}


\par }{\PP \VS{35}Iacov și Ioan, fiii lui Zebedei, s-au apropiat de El și I-au zis: “Învățătorule, vrem să faci pentru noi tot ce-Ți vom cere.”


\par }{\PP \VS{36}El le-a zis:
{\WJ{“Ce vreți să fac pentru voi?”
}}


\par }{\PP \VS{37}Și I-au zis: “Dă-ne nouă să ședem, unul la dreapta Ta și altul la stânga Ta, în slava Ta.”


\par }{\PP \VS{38}Dar Isus le-a zis: “{\WJ{Nu știți ce cereți. Sunteți în stare să beți paharul pe care îl beau Eu și să fiți botezați cu botezul cu care sunt botezat Eu?”
}}


\par }{\PP \VS{39}Ei i-au zis: “Putem.”

\par }{\PP Isus le-a zis: “{\WJ{Voi veți bea paharul pe care îl beau Eu și veți fi botezați cu botezul cu care sunt botezat Eu;
}}

\VS{40}{\WJ{dar șederea la dreapta și la stânga Mea nu este a Mea, ci a cui a fost pregătită.”
}}


\par }{\PP \VS{41}Cei zece, când au auzit, au început să se mânie pe Iacov și pe Ioan.


\par }{\PP \VS{42}Isus i-a chemat și le-a zis: “{\WJ{Știți că cei ce sunt socotiți stăpâni peste neamuri stăpânesc peste ele și că cei mari stăpânesc peste ele.
}}

\VS{43}{\WJ{Dar la voi nu va fi așa, ci oricine vrea să devină mare între voi va fi robul vostru.
}}

\VS{44}{\WJ{Oricine dintre voi va vrea să devină cel dintâi dintre voi va fi robul tuturor.
}}

\VS{45}{\WJ{Căci și Fiul Omului a venit nu ca să fie slujit, ci ca să slujească și să-și dea viața ca răscumpărare pentru mulți.”
}}


\par }{\PP \VS{46}Au ajuns la Ierihon. Pe când ieșea din Ierihon cu ucenicii Săi și cu o mare mulțime, fiul lui Timeu, Bartimeu, un cerșetor orb, ședea lângă drum.

\VS{47}Când a auzit că era Isus Nazarineanul, a început să strige și să spună: “Isuse, Fiul lui David, ai milă de mine!”

\VS{48}Mulți îl mustrau, ca să tacă, dar el striga mult mai tare: “Fiul lui David, ai milă de mine!”


\par }{\PP \VS{49}Isus s-a oprit și a zis:
{\WJ{“Cheamă-l.”
}}

\par }{\PP L-au chemat pe orb și i-au zis: “Înveselește-te! Ridică-te! El te cheamă!”


\par }{\PP \VS{50}El, aruncându-și haina, s-a ridicat și a venit la Isus.


\par }{\PP \VS{51}Isus l-a întrebat:
{\WJ{“Ce vrei să fac pentru tine?”
}}

\par }{\PP Orbul i-a zis: “Rabboni, ca să văd din nou”.


\par }{\PP \VS{52}Isus i-a zis:
{\WJ{“Du-te și mergi. Credința ta te-a făcut sănătos”. Imediat și-a recăpătat vederea și l-a urmat pe Isus pe drum.
}}



\par }\Chap{11}{\PP \VerseOne{1}Când s-au apropiat de Ierusalim, de Betfaghe și Betania, pe Muntele Măslinilor, a trimis doi dintre ucenicii Săi

\VS{2}și le-a zis: “Mergeți
{\WJ{în satul care este în fața voastră. Imediat ce veți intra în el, veți găsi legat un măgăruș tânăr, pe care nu s-a așezat nimeni. Dezlegați-l și aduceți-l.
}}

\VS{3}{\WJ{Dacă vă întreabă cineva: “De ce faceți asta?”, spuneți: “Domnul are nevoie de el”; și îndată îl va trimite înapoi aici.”
}}


\par }{\PP \VS{4}S-au dus și au găsit un măgăruș legat la ușă, afară, în plină stradă, și l-au dezlegat.

\VS{5}Unii dintre cei care stăteau acolo i-au întrebat: “Ce faceți, dezlegând măgărușul?”

\VS{6}Ei le-au spus exact cum le spusese Isus și i-au lăsat să plece.


\par }{\PP \VS{7}Au adus la Isus un măgăruș, și-au aruncat hainele pe el și Isus a șezut pe el.

\VS{8}Mulți își întindeau hainele pe drum, iar alții tăiau ramuri din copaci și le împrăștiau pe drum.

\VS{9}Cei care mergeau în față și cei care veneau după ei strigau: “Osana! Binecuvântat este cel care vine în numele Domnului!

\VS{10}Binecuvântată este împărăția tatălui nostru David, care vine în numele Domnului! Osana în locurile prea înalte!”


\par }{\PP \VS{11}Isus a intrat în Templul din Ierusalim. După ce a privit totul în jur, fiind seară, a ieșit în Betania cu cei doisprezece.


\par }{\PP \VS{12}A doua zi, după ce au ieșit din Betania, i-a fost foame.

\VS{13}Văzând de departe un smochin care avea frunze, a venit să vadă dacă nu cumva ar putea găsi ceva pe el. Când a ajuns la el, nu a găsit decât frunze, căci nu era sezonul smochinelor.

\VS{14}Isus i-a spus:
{\WJ{“Fie ca nimeni să nu mai mănânce vreodată fructe de la tine!”, iar discipolii lui au auzit.
}}


\par }{\PP \VS{15}Au ajuns la Ierusalim, și Isus a intrat în Templu și a început să dea afară pe cei ce vindeau și pe cei ce cumpărau în Templu, să răstoarne mesele schimbătorilor de bani și scaunele celor ce vindeau porumbei.

\VS{16}El nu permitea nimănui să transporte un recipient prin templu.

\VS{17}El îi învăța și le spunea: “{\WJ{Nu este oare scris: “Casa mea va fi numită casă de rugăciune pentru toate neamurile?”? Dar voi ați făcut din ea o peșteră de tâlhari!”
}}


\par }{\PP \VS{18}Preoții cei mai de seamă și cărturarii au auzit acest lucru și căutau cum să-L nimicească. Căci se temeau de el, pentru că toată mulțimea era uimită de învățătura lui.


\par }{\PP \VS{19}Când s-a făcut seară, a ieșit din cetate.

\VS{20}Dimineața, când a trecut pe acolo, a văzut smochinul uscat din rădăcini.

\VS{21}Petru, aducându-și aminte, i-a zis: “Rabi, uite! Smochinul pe care l-ai blestemat s-a uscat”.


\par }{\PP \VS{22}Isus le-a răspuns: “{\WJ{Aveți credință în Dumnezeu.
}}

\VS{23}{\WJ{Căci, cu adevărat vă spun că oricine va spune acestui munte: “Ia-l și aruncă-l în mare” și nu se îndoiește în inima lui, ci crede că ceea ce spune se întâmplă, va avea tot ce spune.
}}

\VS{24}{\WJ{De aceea vă spun că orice lucru pe care îl rugați și îl cereți, credeți că l-ați primit și îl veți avea.
}}

\VS{25}{\WJ{Ori de câte ori stați la rugăciune, iertați, dacă aveți ceva împotriva cuiva, pentru ca și Tatăl vostru, care este în ceruri, să vă ierte greșelile voastre.
}}

\VS{26}{\WJ{Dar dacă nu iertați, nici Tatăl vostru care este în ceruri nu vă va ierta greșelile voastre.”
}}


\par }{\PP \VS{27}Și, întorcându-se la Ierusalim, pe când umbla El în Templu, au venit la El preoții cei mai de seamă, cărturarii și bătrânii

\VS{28}și au început să-I spună: “Cu ce putere faci Tu aceste lucruri? Sau cine ți-a dat această autoritate ca să faci aceste lucruri?”


\par }{\PP \VS{29}Isus le-a zis: “{\WJ{Vă voi pune o întrebare. Răspundeți-mi și vă voi spune cu ce autoritate fac aceste lucruri.
}}

\VS{30}{\WJ{Botezul lui Ioan — a fost din cer sau de la oameni? Răspundeți-mi!”
}}


\par }{\PP \VS{31}Ei se gândeau între ei, zicând: “Dacă vom spune: “Din cer”, el va zice: “De ce nu l-ați crezut?”

\VS{32}Dacă vom spune: “De la oameni””, se temeau de popor, căci toți îl considerau pe Ioan ca fiind cu adevărat un profet.

\VS{33}Ei i-au răspuns lui Isus: “Nu știm”.

\par }{\PP Isus le-a zis:
{\WJ{“Nici Eu nu vă voi spune cu ce putere fac aceste lucruri.”
}}



\par }\Chap{12}{\PP \VerseOne{1}Și a început să le vorbească în parabole.
{\WJ{“Un om a plantat o vie, a pus un gard în jurul ei, a săpat o groapă pentru presa de vin, a construit un turn, a dat-o în arendă unui fermier și a plecat în altă țară.
}}

\VS{2}Când a
{\WJ{venit vremea, a trimis un servitor la fermier pentru a lua de la acesta partea lui din roadele viei.
}}

\VS{3}L-au
{\WJ{prins, l-au bătut și l-au trimis departe cu mâna goală.
}}

\VS{4}{\WJ{El a trimis din nou un alt slujitor la ei; dar ei au aruncat cu pietre în el, l-au rănit la cap și l-au trimis departe cu un tratament rușinos.
}}

\VS{5}{\WJ{A trimis iarăși un altul, și l-au omorât pe el și pe mulți alții, bătându-i pe unii și omorându-i pe alții.
}}

\VS{6}De
{\WJ{aceea, având încă unul, fiul său preaiubit, l-a trimis ultimul la ei, zicând: “Îl vor respecta pe fiul meu”.
}}

\VS{7}{\WJ{Dar țăranii aceia au zis între ei: “Acesta este moștenitorul. Haideți să-l omorâm și moștenirea va fi a noastră'.
}}

\VS{8}L-au
{\WJ{luat, l-au ucis și l-au aruncat afară din vie.
}}

\VS{9}{\WJ{Ce va face, așadar, stăpânul viei? Va veni și îi va nimici pe fermieri și va da via altora.
}}

\VS{10}{\WJ{N-ați citit măcar Scriptura aceasta:
}}

\par }{\Q {\WJ{“Piatra pe care au respins-o zidarii
}}

\par }{\QB {\WJ{a fost numit șef de colț.
}}


\par }{\Q \VS{11}{\WJ{Aceasta a fost de la Domnul.
}}

\par }{\QB {\WJ{Este minunat în ochii noștri”?”?
}}


\par }{\PP \VS{12}Ei au încercat să-L prindă, dar se temeau de mulțime, pentru că au înțeles că vorbea pilda împotriva lor. L-au lăsat și au plecat.

\VS{13}Au trimis la el pe unii dintre farisei și dintre erodieni, ca să-l prindă cu vorba.

\VS{14}După ce au venit, l-au întrebat: “Învățătorule, știm că ești cinstit și că nu te înfrânezi de la nimeni, căci nu ești părtinitor față de nimeni, ci înveți cu adevărat calea lui Dumnezeu. Este sau nu este drept să plătim impozite Cezarului?

\VS{15}Să dăm, sau să nu dăm?”

\par }{\PP Dar El, cunoscând ipocrizia lor, le-a zis:
{\WJ{“De ce Mă puneți la încercare? Aduceți-mi un denar, ca să-l văd”.
}}


\par }{\PP \VS{16}Și au adus-o.

\par }{\PP El le-a întrebat:
{\WJ{“A cui este acest chip și această inscripție?”
}}

\par }{\PP I-au spus: “A lui Cezar”.


\par }{\PP \VS{17}Isus le-a răspuns: “{\WJ{Dați Cezarului ce este al Cezarului și lui Dumnezeu ce este al lui Dumnezeu.”
}}

\par }{\PP Ei s-au minunat foarte mult de el.


\par }{\PP \VS{18}Au venit la El niște saduchei, care ziceau că nu există înviere. L-au întrebat, zicând:

\VS{19}“Învățătorule, Moise ne-a scris: “Dacă un frate al unui om moare și lasă în urma lui o soție și nu lasă copii, fratele său să ia pe soția lui și să ridice urmași fratelui său”.

\VS{20}Erau șapte frați. Primul și-a luat o soție și, murind, nu a lăsat urmași.

\VS{21}Al doilea a luat-o și a murit, fără să lase copii în urma lui. Al treilea, la fel;

\VS{22}iar cei șapte au luat-o și nu au lăsat copii. Ultima dintre toate, femeia a murit și ea.

\VS{23}La înviere, când vor învia, a cui soție va fi ea dintre ei? Căci cei șapte au avut-o ca soție.”


\par }{\PP \VS{24}Isus le-a răspuns: “{\WJ{Nu cumva vă înșelați, pentru că nu cunoașteți Scripturile și puterea lui Dumnezeu?
}}

\VS{25}{\WJ{Căci, când vor învia din morți, nici nu se vor căsători, nici nu vor fi dați în căsătorie, ci vor fi ca îngerii din ceruri.
}}

\VS{26}{\WJ{Dar despre morți, despre faptul că vor învia, n-ați citit în cartea lui Moise, despre Tufișul, cum i-a vorbit Dumnezeu, zicând: “Eu sunt Dumnezeul lui Avraam, Dumnezeul lui Isaac și Dumnezeul lui Iacov”?
}}

\VS{27}{\WJ{El nu este Dumnezeul morților, ci al celor vii. Prin urmare, vă înșelați amarnic”.
}}


\par }{\PP \VS{28}Unul dintre cărturari, venind și auzindu-i pe ei întrebând împreună, și știind că le răspunsese bine, L-a întrebat: “Care poruncă este cea mai mare dintre toate?”


\par }{\PP \VS{29}Isus a răspuns: “{\WJ{Cel mai mare este: “Ascultă, Israele, Domnul Dumnezeul nostru, Domnul este unul.
}}

\VS{30}{\WJ{Să iubești pe Domnul Dumnezeul tău din toată inima ta, din tot sufletul tău, din tot cugetul tău și din toată puterea ta. Aceasta este prima poruncă.
}}

\VS{31}Cea
{\WJ{de-a doua este așa: “Ce este a doua poruncă? “Să iubești pe aproapele tău ca pe tine însuți”. Nu există altă poruncă mai mare decât acestea”.”
}}


\par }{\PP \VS{32}Cărturarul i-a zis: “Cu adevărat, învățătorule, bine ai spus că El este unul și nu este altul decât El;

\VS{33}și a-L iubi pe El cu toată inima, cu toată mintea, cu tot sufletul și cu toată puterea și a iubi pe aproapele ca pe sine însuși este mai de preț decât toate arderile de tot și decât toate jertfele.”


\par }{\PP \VS{34}Isus, văzând că a răspuns cu înțelepciune, i-a zis: “{\WJ{Nu ești departe de Împărăția lui Dumnezeu.”
}}

\par }{\PP Nimeni nu a mai îndrăznit să-i pună întrebări după aceea.

\VS{35}Isus a răspuns, în timp ce învăța în templu:
{\WJ{“Cum se face că scribii spun că Hristosul este fiul lui David?
}}

\VS{36}{\WJ{Căci David însuși a spus în Duhul Sfânt,
}}

\par }{\Q {\WJ{“Domnul a spus Domnului meu,
}}

\par }{\QB {\WJ{“Stai la dreapta mea,
}}

\par }{\QB {\WJ{până ce voi face din vrăjmașii tăi scăunelul picioarelor tale.”
}}


\par }{\PP \VS{37}De
{\WJ{aceea însuși David îl numește Domn, și cum ar putea fi fiul său?”
}}

\par }{\PP Oamenii de rând l-au ascultat cu bucurie.

\VS{38}În învățătura sa, el le spunea: “Feriți-vă
{\WJ{de cărturari, cărora le place să umble în haine lungi, să primească salutări în piețe,
}}

\VS{39}să
{\WJ{obțină cele mai bune locuri în sinagogi și cele mai bune locuri la sărbători,
}}

\VS{40}{\WJ{cei care devorează casele văduvelor și care, pentru a se preface, fac rugăciuni lungi. Aceștia vor primi o condamnare mai mare”.
}}


\par }{\PP \VS{41}Isus a șezut în fața vistieriei și a văzut cum mulțimea arunca bani în vistierie. Mulți dintre cei bogați au aruncat mult.

\VS{42}O văduvă săracă a venit și a aruncat două monede mici de aramă, care sunt cât o monedă de un pătrar.

\VS{43}El și-a chemat discipolii și le-a spus: “{\WJ{Adevărat vă spun că această văduvă săracă a dat mai mult decât toți cei care dau în vistierie,
}}

\VS{44}{\WJ{căcitoți au dat din belșugul lor, dar ea, din sărăcia ei, a dat tot ce avea ca să trăiască.”
}}



\par }\Chap{13}{\PP \VerseOne{1}Când a ieșit din templu, unul din ucenicii Săi I-a zis: “Învățătorule, vezi ce fel de pietre și ce fel de clădiri!”


\par }{\PP \VS{2}Isus i-a zis: “{\WJ{Vezi tu aceste clădiri mari? Nu va rămâne aici piatră pe piatră care să nu fie dărâmată.”
}}


\par }{\PP \VS{3}Pe când ședea El pe Muntele Măslinilor, vizavi de Templu, Petru, Iacov, Ioan și Andrei L-au întrebat în particular:

\VS{4}“Spune-ne, când se vor întâmpla aceste lucruri? Care este semnul că toate aceste lucruri sunt pe cale să se împlinească?”


\par }{\PP \VS{5}Isus, răspunzând, a început să le spună: “Luați
{\WJ{seama ca nimeni să nu vă ducă în rătăcire.
}}

\VS{6}{\WJ{Căci mulți vor veni în numele meu, spunând: “Eu sunt!” și vor duce pe mulți în rătăcire.
}}


\par }{\PP \VS{7}{\WJ{“Când auziți de războaie și de zvonuri de războaie, nu vă tulburați. Căci acestea trebuie să se întâmple, dar sfârșitul nu este încă.
}}

\VS{8}{\WJ{Căci se va ridica națiune împotriva națiunii și regat împotriva regatului. Vor fi cutremure în diferite locuri. Vor fi foamete și necazuri. Aceste lucruri sunt începutul durerilor de naștere.
}}


\par }{\PP \VS{9}{\WJ{Dar păziți-vă, căci vă vor da la consilii. Veți fi bătuți în sinagogi. Veți sta în fața conducătorilor și a împăraților din cauza Mea, ca mărturie pentru ei.
}}

\VS{10}{\WJ{Vestea cea bună trebuie mai întâi să fie propovăduită la toate națiunile.
}}

\VS{11}{\WJ{Când vă vor duce și vă vor preda, nu vă îngrijorați dinainte și nu premeditați ce veți spune, ci spuneți tot ce vi se va da în acel ceas. Căci nu voi sunteți cei care vorbiți, ci Duhul Sfânt.
}}


\par }{\PP \VS{12}“{\WJ{Fratele va da la moarte pe fratele său și tatăl pe copilul său. Copiii se vor ridica împotriva părinților și îi vor face să moară.
}}

\VS{13}Veți
{\WJ{fi urâți de toți oamenii din pricina numelui Meu, dar cel care va răbda până la sfârșit va fi mântuit.
}}


\par }{\PP \VS{14}{\WJ{“Dar când veți vedea urâciunea pustiirii, despre care vorbește proorocul Daniel, stând acolo unde nu trebuie”
}}(să înțeleagă cititorul), “{\WJ{atunci cei din Iudeea să fugă în munți,
}}

\VS{15}iar cel ce
{\WJ{este pe acoperișul casei să nu se coboare și să nu intre să ia nimic din casa lui.
}}

\VS{16}{\WJ{Cel care este pe câmp să nu se întoarcă să își ia mantaua.
}}

\VS{17}{\WJ{Dar vai de cele însărcinate și de cele care alăptează copii în acele zile!
}}

\VS{18}{\WJ{Rugați-vă ca fuga voastră să nu fie în timpul iernii.
}}

\VS{19}{\WJ{Căci în zilele acelea va fi o asuprire, cum nu a mai fost așa ceva de la începutul creației pe care a creat-o Dumnezeu până acum și nici nu va mai fi vreodată.
}}

\VS{20}{\WJ{Dacă Domnul nu ar fi scurtat zilele, nicio făptură nu ar fi fost salvată; dar, de dragul celor aleși, pe care i-a ales, a scurtat zilele.
}}

\VS{21}{\WJ{Așadar, dacă vă spune cineva: “Uite, aici este Hristosul!” sau “Uite, acolo!”, nu-l credeți.
}}

\VS{22}Căci se
{\WJ{vor ridica falși hristoși și falși profeți și vor face semne și minuni, ca să ducă în rătăcire, dacă este posibil, chiar și pe cei aleși.
}}

\VS{23}{\WJ{Dar voi vegheați.
}}

\par }{\PP {\WJ{“Iată, v-am spus toate lucrurile mai dinainte.
}}

\VS{24}{\WJ{Dar în zilele acelea, după acea asuprire, soarele se va întuneca, luna nu-și va mai da lumina,
}}

\VS{25}{\WJ{stelele vor cădea din cer și puterile care sunt în ceruri vor fi zguduite.
}}

\VS{26}{\WJ{Atunci vor vedea pe Fiul Omului venind pe nori, cu mare putere și glorie.
}}

\VS{27}{\WJ{Atunci va trimite pe îngerii Săi și va aduna pe cei aleși ai Săi din cele patru vânturi, de la marginile pământului până la marginile cerului.
}}


\par }{\PP \VS{28}“Și
{\WJ{acum, de la smochin, învățați pilda aceasta. Când acum ramura s-a înmuiat și își produce frunzele, știți că vara este aproape;
}}

\VS{29}tot
{\WJ{așa și voi, când vedeți că se întâmplă aceste lucruri, știți că este aproape, la uși.
}}

\VS{30}{\WJ{Adevărat vă spun că generația aceasta nu va trece până nu se vor întâmpla toate aceste lucruri.
}}

\VS{31}{\WJ{Cerul și pământul vor trece, dar cuvintele Mele nu vor trece.
}}


\par }{\PP \VS{32}{\WJ{Dar despre ziua aceea și despre ceasul acela nimeni nu știe, nici îngerii din ceruri, nici Fiul, ci numai Tatăl.
}}

\VS{33}{\WJ{Vegheați, fiți vigilenți și rugați-vă, căci nu știți când va fi ceasul.
}}


\par }{\PP \VS{34}Este
{\WJ{ca și cum un om care pleacă în altă țară, care a plecat din casa sa și a dat putere slugilor sale, a dat fiecăruia lucrul său și a poruncit portarului să vegheze.
}}

\VS{35}{\WJ{Vegheați, deci, căci nu știți când va veni stăpânul casei — fie seara, fie la miezul nopții, fie când cântă cocoșul, fie dimineața;
}}

\VS{36}{\WJ{ca nu cumva, venind deodată, să vă găsească dormind.
}}

\VS{37}{\WJ{Ceeace vă spun, spun tuturor: Vegheați!”
}}



\par }\Chap{14}{\PP \VerseOne{1}Cu două zile înainte de Paști și de sărbătoarea azimilor, preoții cei mai de seamă și cărturarii căutau cum să-L prindă cu înșelăciune și să-L omoare.

\VS{2}Căci ziceau: “Nu în timpul sărbătorii, pentru că s-ar putea să se facă o răscoală în popor”.


\par }{\PP \VS{3}Pe când era în Betania, în casa lui Simon leprosul, pe când ședea la masă, a venit o femeie cu un vas de alabastru cu mir de nard curat, foarte scump. Ea a spart vasul și l-a turnat pe capul lui.

\VS{4}Dar unii dintre ei erau indignați între ei și ziceau: “De ce s-a risipit mirul acesta?

\VS{5}Pentru că ar fi putut fi vândut cu mai mult de trei sute de denari și dat săracilor.” Așa că au cârtit împotriva ei.


\par }{\PP \VS{6}Dar Isus a zis: “{\WJ{Las-o în pace. De ce o tulburați? Ea a făcut o lucrare bună pentru Mine.
}}

\VS{7}{\WJ{Pentru că voi aveți mereu săraci cu voi și, oricând vreți, le puteți face bine, dar pe mine nu mă veți avea mereu.
}}

\VS{8}{\WJ{Ea a făcut tot ce a putut. Mi-a uns trupul mai dinainte pentru înmormântare.
}}

\VS{9}Cu
{\WJ{siguranță vă spun că, oriunde se va propovădui această Bună Vestire în toată lumea, se va vorbi și despre ceea ce a făcut această femeie, ca o amintire a ei.”
}}


\par }{\PP \VS{10}Iuda Iscarioteanul, care era unul din cei doisprezece, s-a dus la preoții cei mai de seamă, ca să li-l predea.

\VS{11}Aceștia, când au auzit, s-au bucurat și au făgăduit să-i dea bani. El a căutat cum ar putea să-l livreze în mod convenabil.


\par }{\PP \VS{12}În prima zi a azimilor, când se jertfeau Paștile, ucenicii Lui L-au întrebat: “Unde vrei să mergem să pregătim ca să mâncați Paștile?”


\par }{\PP \VS{13}Și a trimis doi dintre ucenicii Săi și le-a zis: “{\WJ{Mergeți în cetate și vă va întâmpina un om care duce un ulcior cu apă. Urmați-l
}}

\VS{14}și,
{\WJ{oriunde va intra, spuneți-i stăpânului casei: “Învățătorul spune: “Unde este camera de oaspeți, unde pot mânca Paștele cu discipolii mei?””.
}}

\VS{15}{\WJ{El însuși vă va arăta o cameră mare de sus, mobilată și pregătită. Pregătește-ne acolo”.
}}


\par }{\PP \VS{16}Ucenicii Lui au ieșit, au intrat în cetate și au găsit lucrurile așa cum le spusese El și au pregătit Paștele.


\par }{\PP \VS{17}Când s-a făcut seară, a venit cu cei doisprezece.

\VS{18}Pe când ședeau și mâncau, Isus a zis: “{\WJ{Vă spun cu siguranță că unul dintre voi Mă va trăda — cel care mănâncă cu Mine.”
}}


\par }{\PP \VS{19}Și au început să se întristeze și să-L întrebe unul câte unul: “Nu cumva sunt eu?” Și un altul a zis: “Cu siguranță nu eu?”


\par }{\PP \VS{20}El le-a răspuns: “{\WJ{Este unul din cei doisprezece, cel care se scufundă cu Mine în mâncare.
}}

\VS{21}{\WJ{Căci Fiul Omului merge așa cum este scris despre El, dar vai de omul acela prin care este trădat Fiul Omului! Mai bine pentru acel om ar fi fost mai bine să nu se fi născut.”
}}


\par }{\PP \VS{22}Pe când mâncau ei, Isus a luat o pâine și, după ce a binecuvântat-o, a frânt-o, le-a dat-o și le-a zis: “{\WJ{Luați, mâncați. Acesta este trupul Meu”.
}}


\par }{\PP \VS{23}A luat paharul și, după ce a mulțumit, le-a dat. Toți au băut din el.

\VS{24}Și le-a zis: “{\WJ{Acesta este sângele Meu, al noului legământ, care se varsă pentru mulți.
}}

\VS{25}{\WJ{Adevărat vă spun că nu voi mai bea din rodul viței de vie până în ziua când îl voi bea din nou în Împărăția lui Dumnezeu.”
}}

\VS{26}După ce au cântat un imn, au ieșit pe Muntele Măslinilor.


\par }{\PP \VS{27}Isus le-a zis: “În
{\WJ{noaptea aceasta, toți vă veți poticni din pricina Mea, căci este scris: “Voi lovi pe păstor și oile vor fi risipite”.
}}

\VS{28}{\WJ{Totuși, după ce voi fi înviat, voi merge înaintea voastră în Galileea.”
}}


\par }{\PP \VS{29}Dar Petru i-a zis: “Chiar dacă toți se vor supăra, eu nu mă voi supăra.”


\par }{\PP \VS{30}Isus i-a zis: “{\WJ{Adevărat îți spun că astăzi, chiar în noaptea aceasta, înainte de a cânta cocoșul de două ori, te vei lepăda de Mine de trei ori.”
}}


\par }{\PP \VS{31}Dar el a vorbit cu atât mai mult: “Chiar dacă va trebui să mor cu voi, nu mă voi lepăda de voi.” Și toți au spus același lucru.


\par }{\PP \VS{32}Și au ajuns la un loc numit Ghetsimani. El a zis ucenicilor Săi:
{\WJ{“Ședeți aici, în timp ce Mă voi ruga.”
}}

\VS{33}A luat cu El pe Petru, Iacov și Ioan și a început să fie foarte tulburat și necăjit.

\VS{34}El le-a zis: “{\WJ{Sufletul Meu este foarte întristat, până la moarte. Rămâneți aici și vegheați”.
}}


\par }{\PP \VS{35}A înaintat puțin, a căzut la pământ și s-a rugat ca, dacă se putea, să treacă de la el ceasul acela.

\VS{36}Și a zis: “{\WJ{Abba, Tată, toate lucrurile sunt cu putință pentru Tine. Te rog, îndepărtează de la mine paharul acesta. Oricum, nu ceea ce doresc eu, ci ceea ce dorești Tu”.
}}


\par }{\PP \VS{37}El a venit, i-a găsit dormind și a zis lui Petru:
{\WJ{“Simon, dormi? N-ați putut să vegheați un ceas?
}}

\VS{38}{\WJ{Vegheați și rugați-vă, ca să nu intrați în ispită. Într-adevăr, duhul este binevoitor, dar carnea este slabă.”
}}


\par }{\PP \VS{39}Și s-a dus iarăși și s-a rugat, zicând aceleași cuvinte.

\VS{40}S-a întors iarăși și i-a găsit dormind, căci aveau ochii foarte îngreunați; și nu știau ce să-i răspundă.

\VS{41}A venit a treia oară și le-a zis:
{\WJ{“Dormiți acum și odihniți-vă. Este de ajuns. A sosit ceasul. Iată, Fiul Omului este dat în mâinile păcătoșilor.
}}

\VS{42}{\WJ{Ridicați-vă! Haideți să mergem. Iată, cel care mă trădează este aproape”.
}}


\par }{\PP \VS{43}Și îndată, pe când vorbea El încă, a venit Iuda, unul din cei doisprezece, și cu el o mulțime de oameni cu săbii și cu ciomege, de la preoții cei mai de seamă, de la cărturari și de la bătrâni.

\VS{44}Cel care l-a trădat le dăduse un semn, zicând: “Pe cine voi săruta eu, acela este. Prindeți-l și duceți-l în siguranță”.

\VS{45}După ce a venit, îndată s-a apropiat de el și i-a zis: “Rabi! Rabi!” și l-a sărutat.

\VS{46}Ei au pus mâinile pe el și l-au prins.

\VS{47}Dar unul dintre cei care stăteau acolo a scos sabia și a lovit pe slujitorul marelui preot și i-a tăiat urechea.


\par }{\PP \VS{48}Isus le-a răspuns: “Ați
{\WJ{ieșit voi, ca un tâlhar, cu săbii și cu ciomege, ca să Mă prindeți?
}}

\VS{49}{\WJ{Eu eram zilnic cu voi în templu, învățând, și nu m-ați arestat. Dar aceasta s-a întâmplat ca să se împlinească Scripturile”.
}}


\par }{\PP \VS{50}Toți l-au lăsat și au fugit.

\VS{51}Și un tânăr l-a urmat, cu o pânză de in înfășurată peste trupul lui gol. Tinerii l-au apucat,

\VS{52}dar el a lăsat pânza de in și a fugit de ei gol.

\VS{53}L-au dus pe Isus la marele preot. Toți preoții cei mai de seamă, bătrânii și cărturarii au venit împreună cu el.


\par }{\PP \VS{54}Petru Îl urmase de departe, până ce a intrat în curtea marelui preot. El stătea cu ofițerii și se încălzea la lumina focului.

\VS{55}Preoții cei mai de seamă și tot consiliul căutau martori împotriva lui Isus, ca să-L condamne la moarte, și n-au găsit niciunul.

\VS{56}Căci mulți dădeau mărturie falsă împotriva lui și mărturiile lor nu se potriveau între ele.

\VS{57}Unii s-au ridicat în picioare și au depus mărturie mincinoasă împotriva lui, zicând:

\VS{58}“L-am auzit spunând: “Voi nimici templul acesta făcut de mână și în trei zile voi zidi un altul făcut fără mâini”.”

\VS{59}Chiar și așa, mărturia lor nu era de acord.


\par }{\PP \VS{60}Marele preot s-a ridicat în picioare la mijloc și a întrebat pe Isus: “Nu ai răspuns? Ce este ceea ce mărturisesc aceștia împotriva ta?”.

\VS{61}Dar el a tăcut și nu a răspuns nimic. Din nou, marele preot l-a întrebat: “Tu ești Hristosul, Fiul Celui Binecuvântat?”.


\par }{\PP \VS{62}Isus a spus:
{\WJ{“Eu sunt. Îl veți vedea pe Fiul Omului șezând la dreapta Puterii și venind cu norii cerului.”
}}


\par }{\PP \VS{63}Marele preot și-a rupt hainele și a zis: “Ce nevoie mai avem de martori?

\VS{64}Ați auzit blasfemia! Ce părere aveți?” Cu toții l-au condamnat ca fiind vrednic de moarte.

\VS{65}Unii au început să-l scuipe, să-i acopere fața, să-l bată cu pumnii și să-i spună: “Profetizează!” Ofițerii l-au lovit cu palmele.


\par }{\PP \VS{66}Pe când Petru era în curtea de jos, a venit una din slujnicele marelui preot,

\VS{67}și, văzând pe Petru că se încălzea, s-a uitat la el și i-a zis: “Și tu ai fost cu Nazarineanul Iisus!”.


\par }{\PP \VS{68}Dar el a negat, zicând: “Nu știu și nu înțeleg ce spui.” A ieșit pe verandă, iar cocoșul a cântat.


\par }{\PP \VS{69}Și servitoarea, văzându-l, a început iarăși să spună celor ce stăteau de față: “Acesta este unul dintre ei”.

\VS{70}Dar el a negat din nou. După puțină vreme, cei care stăteau de față i-au spus din nou lui Petru: “Cu adevărat, tu ești unul dintre ei, căci ești galileean și vorbirea ta o arată.”

\VS{71}Dar el a început să blesteme și să jure: “Nu-l cunosc pe omul acesta despre care vorbiți!”


\par }{\PP \VS{72}Cocoșul a cântat a doua oară. Petru și-a amintit cuvintele pe care i le spusese Isus:
{\WJ{“Înainte ca cocoșul să cânte de două ori, te vei lepăda de Mine de trei ori”. Când s-a gândit la asta, a plâns.
}}



\par }\Chap{15}{\PP \VerseOne{1}Îndată de dimineață, preoții cei mai de seamă, împreună cu bătrânii, cărturarii și tot consiliul, s-au sfătuit, au legat pe Isus, L-au dus și L-au predat lui Pilat.

\VS{2}Pilat l-a întrebat: “Ești Tu regele iudeilor?”.

\par }{\PP El a răspuns:
{\WJ{“Așa spui tu”.
}}


\par }{\PP \VS{3}Preoții cei mai de seamă L-au acuzat de multe lucruri.

\VS{4}Pilat l-a întrebat din nou: “Nu ai răspuns? Vezi câte lucruri depun împotriva ta!”


\par }{\PP \VS{5}Dar Isus n-a mai răspuns nimic, așa că Pilat s-a mirat.


\par }{\PP \VS{6}Și la sărbătoare le dădea drumul la un prizonier, pe oricine îi cereau.

\VS{7}Era unul numit Baraba, legat împreună cu tovarășii săi de insurecție, oameni care, în timpul insurecției, săvârșiseră o crimă.

\VS{8}Mulțimea, strigând cu voce tare, a început să-i ceară să facă pentru el ceea ce făcea întotdeauna pentru ei.

\VS{9}Pilat le-a răspuns: “Vreți să vi-l eliberez pe Regele iudeilor?”

\VS{10}Pentru că a înțeles că preoții de seamă l-au predat din invidie.

\VS{11}Dar preoții cei mai de seamă au stârnit mulțimea, ca el să le elibereze în schimb pe Baraba.

\VS{12}Pilat i-a întrebat din nou: “Ce să fac atunci cu cel pe care voi îl numiți Regele iudeilor?”


\par }{\PP \VS{13}Și au strigat din nou: “Răstignește-L!”


\par }{\PP \VS{14}Pilat le-a zis: “Ce rău a făcut?”

\par }{\PP Dar ei au strigat foarte tare: “Răstignește-L!”


\par }{\PP \VS{15}Pilat, vrând să mulțumească norodul, i-a dat drumul lui Baraba, iar pe Isus, după ce l-a biciuit, l-a dat spre răstignire.


\par }{\PP \VS{16}Soldații l-au dus în curtea care este pretoriul și au chemat toată cohorta.

\VS{17}L-au îmbrăcat cu purpură; și i-au împletit o coroană de spini, pe care i-au pus-o.

\VS{18}Au început să-l salute: “Slavă, rege al iudeilor!”.

\VS{19}I-au lovit capul cu o trestie, l-au scuipat și, plecându-și genunchii, i-au adus omagiu.

\VS{20}După ce și-au bătut joc de el, i-au luat mantia de purpură și i-au pus hainele lui. L-au dus afară ca să-l răstignească.


\par }{\PP \VS{21}Și au silit pe unul care trecea pe acolo, venind de la țară, Simon din Cirene, tatăl lui Alexandru și al lui Rufus, să meargă cu ei, ca să-și poarte crucea.

\VS{22}L-au dus la locul numit Golgota, care înseamnă, fiind interpretat, “Locul unui craniu”.

\VS{23}I-au oferit să bea vin amestecat cu smirnă, dar el nu l-a luat.


\par }{\PP \VS{24}L-au răstignit, I-au împărțit hainele Lui și au tras la sorți ce să ia fiecare.

\VS{25}Era ceasul al treilea când L-au răstignit.

\VS{26}Și peste el era scrisă supranumele acuzării: “REGELE IUDEILOR”.

\VS{27}Împreună cu El au răstignit doi tâlhari, unul la dreapta și altul la stânga.

\VS{28}S-a împlinit Scriptura care zice: “A fost socotit împreună cu călcătorii de lege”.


\par }{\PP \VS{29}Cei ce treceau pe acolo Îl huleau, clătinându-și capetele și zicând: “Tu, care ai distrus templul și l-ai zidit în trei zile,

\VS{30}mântuiește-te și coboară-te de pe cruce!”


\par }{\PP \VS{31}La fel și preoții cei mai de seamă, care își băteau joc între ei împreună cu cărturarii, ziceau: “A mântuit pe alții. El nu se poate salva pe sine însuși.

\VS{32}Hristosul, Împăratul lui Israel, să se coboare acum de pe cruce, ca să-l vedem și să-l credem.” Cei care au fost răstigniți împreună cu el l-au insultat și ei.


\par }{\PP \VS{33}Când a venit ceasul al șaselea, a fost întuneric peste toată țara până la ceasul al nouălea.

\VS{34}La ceasul al nouălea, Isus a strigat cu glas tare, zicând:
{\WJ{“Eloi, Eloi, lama sabachthani?”}}, ceea ce înseamnă, fiind interpretat:
{\WJ{“Dumnezeul Meu, Dumnezeul Meu, pentru ce M-ai părăsit?”}}


\par }{\PP \VS{35}Unii din cei ce stăteau acolo, auzind, au zis: “Iată, cheamă pe Ilie.”


\par }{\PP \VS{36}Unul a alergat și, umplând un burete plin de oțet, l-a pus pe o trestie și i-a dat să bea, zicând: “Lasă-l în pace. Să vedem dacă vine Ilie să îl dea jos”.


\par }{\PP \VS{37}Isus a strigat cu glas tare și și-a dat duhul.

\VS{38}Voalul templului s-a rupt în două, de sus până jos.

\VS{39}Când centurionul, care stătea vizavi de el, a văzut că a strigat așa și că și-a dat ultima suflare, a zis: “Cu adevărat, omul acesta era Fiul lui Dumnezeu!”


\par }{\PP \VS{40}Erau și femei care priveau de departe, printre care Maria Magdalena, Maria, mama lui Iacov cel mic și a lui Iose și Salomeea,

\VS{41}care, când era în Galileea, L-au urmat și Îi slujeau, și multe alte femei care s-au suit cu El la Ierusalim.


\par }{\PP \VS{42}Și când s-a făcut seară, fiindcă era ziua pregătirii, adică ziua dinaintea Sabatului,

\VS{43}a venit Iosif din Arimateea, un sfetnic de seamă, care și el însuși aștepta Împărăția lui Dumnezeu. A intrat cu îndrăzneală la Pilat și a cerut trupul lui Isus.

\VS{44}Pilat a fost surprins să audă că era deja mort; și chemând centurionul, l-a întrebat dacă era mort de mult timp.

\VS{45}Când a aflat de la centurion, i-a dat trupul lui Iosif.

\VS{46}Acesta a cumpărat o pânză de in și, coborându-l, l-a înfășurat în pânză de in și l-a pus într-un mormânt care fusese tăiat într-o stâncă. A rostogolit o piatră la ușa mormântului.

\VS{47}Maria Magdalena și Maria, mama lui Iose, au văzut unde fusese pus.



\par }\Chap{16}{\PP \VerseOne{1}După ce a trecut Sabatul, Maria Magdalena, Maria, mama lui Iacov, și Salomeea au cumpărat miresme, ca să vină să-L ungă.

\VS{2}Foarte devreme, în prima zi a săptămânii, au venit la mormânt, pe când răsărise soarele.

\VS{3}Și ziceau între ele: “Cine va rostogoli pentru noi piatra de la ușa mormântului?”

\VS{4}căci era foarte mare. Uitându-se în sus, au văzut că piatra era rostogolită.


\par }{\PP \VS{5}Intrând în mormânt, au văzut un tânăr care ședea în partea dreaptă, îmbrăcat într-o haină albă, și au rămas înmărmuriți.

\VS{6}El le-a zis: “Nu vă mirați. Voi îl căutați pe Isus, Nazarineanul, care a fost răstignit. El a înviat! El nu este aici. Vedeți locul unde l-au pus!

\VS{7}Dar mergeți și spuneți-le discipolilor săi și lui Petru: “El merge înaintea voastră în Galileea. Acolo îl veți vedea, așa cum v-a spus””.”


\par }{\PP \VS{8}Ei au ieșit și au fugit de la mormânt, pentru că tremurau și erau înmărmuriți. Nu au spus nimic nimănui, căci le era frică.


\par }{\PP \VS{9}În prima zi a săptămânii, când S-a sculat dis-de-dimineață, S-a arătat mai întâi Mariei Magdalena, din care scosese șapte demoni.

\VS{10}Ea s-a dus și a anunțat pe cei care fuseseră cu El, în timp ce se jeleau și plângeau.

\VS{11}Când au auzit că era viu și că fusese văzut de ea, nu au crezut.


\par }{\PP \VS{12}După aceea, El s-a arătat sub o altă formă, la doi dintre ei, pe când mergeau pe jos, în drum spre țară.

\VS{13}Ei s-au dus și au povestit celorlalți. Nici ei nu i-au crezut.


\par }{\PP \VS{14}După aceea s-a arătat celor unsprezece, pe când ședeau la masă, și i-a mustrat pentru necredința și împietrirea inimii lor, pentru că nu crezuseră pe cei ce-l văzuseră după ce înviase.

\VS{15}El le-a spus:
{\WJ{“Mergeți în toată lumea și propovăduiți Vestea cea Bună la toată creația.
}}

\VS{16}Cine va
{\WJ{crede și va fi botezat va fi mântuit, dar cine nu va crede va fi condamnat.
}}

\VS{17}{\WJ{Aceste semne îi vor însoți pe cei care vor crede: în numele meu vor alunga demoni, vor vorbi limbi noi,
}}

\VS{18}vor lua în
{\WJ{brațe șerpi și, dacă vor bea ceva mortal, nu le va face niciun rău, își vor pune mâinile peste bolnavi și aceștia se vor însănătoși.”
}}


\par }{\PP \VS{19}Domnul,
\FTNT{*}{{\FR 16:19 }{\FT{NU adaugă “Isus”}}}după ce le-a vorbit, S-a înălțat la cer și a șezut la dreapta lui Dumnezeu.

\VS{20}Ei au ieșit și propovăduiau pretutindeni, Domnul lucrând cu ei și confirmând cuvântul prin semnele care urmau. Amin.


\par }